\documentclass[letterpaper,twocolumn,10pt]{article}
\usepackage{usenix-2020-09}
\usepackage{amsmath,amssymb,mathtools,amsthm}
\usepackage{booktabs,multirow,graphicx,framed}
\usepackage[advantage,adversary,notions,probability,events,complexity]{cryptocode}

\emergencystretch=1.5em

\newtheorem{theorem}{Theorem}
\newtheorem{lemma}[theorem]{Lemma}
\newtheorem{corollary}[theorem]{Corollary}
\theoremstyle{definition}
\newtheorem{definition}[theorem]{Definition}

\newcommand{\F}{\mathbb{F}}\newcommand{\B}{\mathbb{B}}
\newcommand{\Z}{\mathbb{Z}}\newcommand{\N}{\mathbb{N}}
\newcommand{\U}{\mathcal{U}}\newcommand{\cF}{\mathcal{F}}
\newcommand{\cL}{\mathcal{L}}\newcommand{\cS}{\mathcal{S}}
\newcommand{\cA}{\mathcal{A}}\newcommand{\cB}{\mathcal{B}}\newcommand{\cZ}{\mathcal{Z}}
\newcommand{\prm}{\mathsf{prm}}\newcommand{\Enc}{\mathsf{Enc}}
\newcommand{\Dec}{\mathsf{Dec}}\newcommand{\Encode}{\mathsf{Encode}}
\newcommand{\DeleteSet}{\mathsf{DeleteSet}}\newcommand{\Peel}{\mathsf{Peel}}
\newcommand{\List}{\mathsf{List}}\newcommand{\uPeel}{\mathsf{uPeel}}
\newcommand{\PRF}{\mathsf{PRF}}\newcommand{\ChaCha}{\mathsf{ChaCha20}}
\newcommand{\Adv}{\mathsf{Adv}}\newcommand{\negl}{\mathsf{negl}}\newcommand{\poly}{\mathsf{poly}}
\newcommand{\sample}{\leftarrow_{\$}}\newcommand{\concat}{\|}

\begin{document}
\title{Forward-Private Updatable Private Set Union from Incremental IBLT Encoding}
\author{Anonymous Submission}
\maketitle

% ============================================================
\begin{abstract}
Private Set Union (PSU) enables two parties to compute the union of their
private sets. While recent IBLT-based protocols achieve PSU performance
competitive with Private Set Intersection, they are \emph{static}: computing the
union on updated sets requires re-running the full protocol from scratch, with
cost proportional to total set size rather than the size of the update.

We present \textbf{F-PSU}, the first forward-private updatable PSU protocol.
F-PSU exploits the linearity of IBLT encoding to support incremental
(delta-based) set updates, achieving per-epoch communication and computation
$O(|\Delta|)$ with overwhelming probability. Forward privacy is obtained via a
PRF-based key evolution chain with immediate erasure, adding negligible overhead.
We also prove a \emph{leakage lower bound}: any updatable PSU protocol with
sublinear communication must leak individual add and delete counts---our
protocol matches this bound and is therefore leakage-optimal.

We further show that a single execution of the UnionPeel sub-protocol
simultaneously recovers the union, the intersection, and their
cardinalities---yielding a \emph{unified} updatable private set operation
framework (updPSO) with zero additional communication or computation.

We implement F-PSU in C++ and evaluate it across set sizes $2^{10}$ to $2^{20}$
under four network profiles (LAN to 5~Mbps WAN). For a $0.1\%$ update ratio,
per-epoch communication is $643\times$ smaller than the static baseline.
Incremental peeling examines only $O(k|\Delta|)$ cells, with cascade never
exceeding $2\times k|\Delta|$ and $100\%$ recovery rate. Amortized over 200
epochs, IBLT operations are reduced by $114\times$. All code is open source.
\end{abstract}

% ============================================================
\section{Introduction}
% ============================================================

Private Set Union (PSU)~\cite{KRT19,JLZ20,ZCL+23,Tu+25} allows two parties to
compute the union $X \cup Y$ of their private sets while revealing nothing
beyond the union itself. With applications in privacy-preserving identity
management, collaborative cyber-risk assessment, private database joins, and
federated learning~\cite{DCW13,IKN+20,KRS+19}, PSU is a fundamental secure
computation primitive alongside Private Set Intersection (PSI).

Historically, PSU has been significantly slower than PSI---often by $30\times$
or more~\cite{PT26,KRT19}. The recent breakthrough of Piske and
Trieu~\cite{PT26} (Eurocrypt 2026) closed this gap by introducing the first
IBLT-based PSU protocol. Their key insight, \emph{union peelability}, enables
peeling the union of two IBLTs without materializing the combined table, using
oblivious transfer (OT) and an oblivious PRF (OPRF). The resulting protocol
computes PSU on $2^{20}$ elements in under 3 seconds, matching state-of-the-art
PSI performance.

\textbf{The static limitation.} Piske and Trieu's protocol, like all prior PSU
constructions, is \emph{static}: both parties must encode their entire sets into
fresh IBLTs for each computation. In many real-world applications, however,
sets evolve incrementally---new elements are added, old ones expire. Consider a
collaborative threat-intelligence platform where multiple organizations share
indicators of compromise (IoCs): hundreds of new IoCs appear daily, old ones are
retired, and the aggregated union must be recomputed to maintain situational
awareness. Re-running the full protocol on each update wastes resources
proportional to the total set size (potentially millions of IoCs), even when
only a handful of elements have changed. Similar dynamics arise in
privacy-preserving contact discovery, federated learning with evolving client
datasets, and private database join maintenance under incremental updates.

\textbf{Why not batch?} Batching updates and re-running static PSU periodically
reduces computational overhead but creates a window where the computed union is
stale. In threat intelligence, a 24-hour stale union means missing active
attack indicators. Updatable PSU eliminates this tension by making per-epoch
cost proportional to the update size, enabling frequent recomputation without
the penalty of full re-encoding.

\textbf{The forward privacy requirement.} The multi-epoch nature of updatable
protocols introduces a new security concern: \emph{forward privacy}. If a party
is compromised at epoch $\tau$, the adversary should not be able to recover the
party's historical set elements from epochs $t < \tau$, even having observed all
protocol messages. This is the privacy analogue of forward secrecy in
Signal~\cite{Signal}, and is critical when sets contain sensitive data with
limited lifetimes. Privacy regulations (GDPR, CCPA) mandate data retention
limits; if historical elements remain decryptable after key compromise, the
system violates the principle that deleted data should be irrecoverable.

\textbf{Existing updatable work.} Several updatable PSI protocols
exist~\cite{BMX22,BMS+24,LMR+26,ABF+25,LTS+26}, with FUPSI~\cite{WZC+24}
achieving forward privacy via keyword PIR. The enhanced PSU of Tu et
al.~\cite{Tu+25} eliminates within-execution leakage but does not address
cross-epoch forward privacy. \textbf{No prior work} addresses updatable PSU,
let alone forward-private updatable PSU.

\textbf{Our contributions.} We initiate the study of updatable PSU with forward
privacy:

\begin{enumerate}
\item \textbf{The F-PSU protocol} (Section~\ref{sec:protocol}). The first
forward-private updatable PSU protocol, supporting \emph{both balanced and
unbalanced} settings natively. Built on three mechanisms: (i)~delta IBLT
encoding for $O(k|\Delta|)$ maintenance; (ii)~split-path incremental peeling
with a provable cascade bound; (iii)~PRF-based key evolution with immediate
erasure, providing forward privacy with zero additional communication.
A natural byproduct of the UnionPeel sub-protocol is that it simultaneously
recovers the \emph{intersection} and cardinalities at zero extra
cost---extending F-PSU to a unified updatable private set operation
framework (updPSO).

\item \textbf{Leakage lower bound} (Section~3, Theorem~\ref{thm:lower}). We prove that
\emph{any} $1$-updatable PSU protocol must leak the individual add and delete
counts for each party---an inherent efficiency-leakage trade-off. F-PSU leaks
exactly these counts, making it \textbf{leakage-optimal}. The proof relies
only on the semi-honest model and Shannon's source coding theorem.

\item \textbf{Adaptive security} (Section~\ref{sec:security}). We formalize
the notion of adaptive input selection for updatable PSU---where an adversary
chooses deltas based on prior transcripts---motivated by recent attacks on
UPSI~\cite{LMR+26}. We prove that F-PSU is \emph{naturally immune} to this
attack class with zero protocol modification, as a structural consequence of
PSU's superset output. This contrasts sharply with PSI, where explicit
countermeasures are required.

\item \textbf{Formal security} (Section~\ref{sec:security}). We prove F-PSU
UC-realizes~\cite{Can01} $\cF_{\mathsf{updPSO}}$ in the $(\cF_{\mathsf{OPRF}},
\cF_{\mathsf{OT}})$-hybrid model against semi-honest adaptive adversaries,
for all four output modes (PSU, PSI, PSU-CA, PSI-CA). The modular proof
separates per-epoch and cross-epoch security, enabling independent upgrades.
A standalone game-based definition (FP-IND) cleanly captures the PRF-chain
guarantee.

\item \textbf{Implementation and evaluation} (Section~\ref{sec:eval}). C++
implementation ($\approx$2,000 lines) using libOTe~\cite{libOTe}, evaluated
across $2^{10}$ to $2^{20}$ elements under four network profiles (LAN to
5~Mbps WAN). At $0.1\%$ update ratio: $643\times$ communication reduction and
$114\times$ amortized IBLT operation reduction. A self-healing key chain
extension with zero additional communication is in
Appendix~\ref{sec:selfheal}.
\end{enumerate}

% ============================================================
\section{Preliminaries}
\label{sec:prelim}
% ============================================================

\subsection{Invertible Bloom Lookup Tables}

Invertible Bloom Lookup Tables (IBLTs) were introduced by Goodrich and
Mitzenmacher~\cite{GM11}. An IBLT is a variant of Bloom filter that supports,
in addition to inserting and deleting elements, the ability to \emph{list} all
stored elements. Unlike a standard Bloom filter---which only answers membership
queries with a configurable false-positive rate and requires $\approx 60\times$
expansion for negligible error---an IBLT stores key-value pairs in its cells
and achieves listing with expansion factor $1.5$--$4.5\times$ (for $k = 4$ hash
functions). This makes IBLTs uniquely suited for set reconciliation and private
set union, where the goal is to recover the entire set rather than answer
individual membership queries.

An IBLT is parameterized by $\prm = (M, k, \ell, H)$: a modulus $M \in
\N_{\geq 2}$ (typically $2^{64}$), $k$ hash functions $H = \{h_i : \U \to
[\ell]\}_{i \in [k]}$ each mapping to a distinct subtable of size $\ell$, and
two $k \times \ell$ matrices $\mathsf{sum}$ (values in $\Z_M$) and
$\mathsf{cnt}$ (non-negative integers). We write $m = k\ell$ for the total
number of cells.

The three core operations are given in Algorithms~\ref{alg:encode}--\ref{alg:list}.

\begin{figure}[h]
\fbox{\begin{minipage}{0.95\columnwidth}
\smallskip
\noindent\textbf{Algorithm 1:} $\Encode_{\prm}(X)$ --- IBLT encoding.
\label{alg:encode}

\smallskip\noindent\textit{Input:} Set $X \subseteq \U$, empty IBLT $(\mathsf{sum}, \mathsf{cnt})$

\noindent\textit{Output:} IBLT containing all $x \in X$
\begin{enumerate}
\item For each $x \in X$, $i \in [k]$:
  \begin{enumerate}
  \item $j \leftarrow h_i(x)$
  \item $\mathsf{sum}[i][j] \leftarrow \mathsf{sum}[i][j] + x \pmod M$
  \item $\mathsf{cnt}[i][j] \leftarrow \mathsf{cnt}[i][j] + 1$
  \end{enumerate}
\end{enumerate}
\smallskip
\end{minipage}}
\end{figure}

\begin{figure}[h]
\fbox{\begin{minipage}{0.95\columnwidth}
\smallskip
\noindent\textbf{Algorithm 2:} $\DeleteSet_{\prm}(\mathsf{IBLT}, X)$ --- batch deletion.
\label{alg:deleteset}

\smallskip\noindent\textit{Input:} IBLT, set $X$ to delete

\noindent\textit{Effect:} Removes all $x \in X$ from the IBLT
\begin{enumerate}
\item For each $x \in X$, $i \in [k]$:
  \begin{enumerate}
  \item $j \leftarrow h_i(x)$
  \item $\mathsf{sum}[i][j] \leftarrow \mathsf{sum}[i][j] - x \pmod M$
  \item $\mathsf{cnt}[i][j] \leftarrow \mathsf{cnt}[i][j] - 1$
  \end{enumerate}
\end{enumerate}
\smallskip
\end{minipage}}
\end{figure}

IBLT encoding is \emph{linear}: $\Encode(A \cup B) = \Encode(A) \oplus
\Encode(B)$ and $\Encode(A \setminus B) = \Encode(A) \ominus \Encode(B)$,
where $\oplus$/$\ominus$ denote cell-wise modular addition/subtraction.

A cell with $\mathsf{cnt}[i][j] = 1$ is a \emph{singleton}: its sum equals
the unique element stored there.

\begin{figure}[h]
\fbox{\begin{minipage}{0.95\columnwidth}
\smallskip
\noindent\textbf{Algorithm 3:} $\List(\mathsf{IBLT})$ --- full listing (peeling).
\label{alg:list}

\smallskip\noindent\textit{Input:} IBLT with parameters $\prm$

\noindent\textit{Output:} Set of all encoded elements, or $\bot$ on failure
\begin{enumerate}
\item While there exists a singleton cell $(i,j)$ with $\mathsf{cnt}[i][j] = 1$:
  \begin{enumerate}
  \item $x \leftarrow \mathsf{sum}[i][j]$ \hfill // recover unique element
  \item Output $x$
  \item $\DeleteSet(\mathsf{IBLT}, \{x\})$ \hfill // remove from all $k$ cells
  \end{enumerate}
\item If no non-empty cells remain, output all recovered elements; otherwise fail.
\end{enumerate}
\smallskip
\end{minipage}}
\end{figure}

\begin{theorem}[IBLT Listing~\cite{GM11,PT26}]\label{thm:iblt}
For $m = k\ell > (c_k + \varepsilon)\tau$ with $n \leq \tau$ elements,
$\Pr[\List\text{ fails}] = O(\tau^{-k+2})$. With $k = 4$, expansion
$m/\tau \in [1.5, 4.5]$, failure $\leq 2^{-\lambda}$.
\end{theorem}

In the above, $c_k$ is the threshold constant for $2$-core emergence ($c_3
\approx 1.222$, $c_4 \approx 1.295$). If $\List$ terminates with non-singleton
cells remaining (a $2$-core), listing has failed. The expansion factor $m/n$
must exceed $c_k$ to guarantee an empty $2$-core with overwhelming probability.

\subsection{Union Peelability}

Piske and Trieu~\cite{PT26} observed that the union of two IBLTs can be peeled
without materializing the combined table. For $\mathsf{IBLT}_0 =
\Encode(X_0)$ and $\mathsf{IBLT}_1 = \Encode(X_1)$, the virtual union is
$\mathsf{IBLT}_\cup = \mathsf{IBLT}_0 \oplus \mathsf{IBLT}_1$.

\begin{definition}[$\uPeel$~\cite{PT26}]\label{def:upeel}
\footnotesize
\[
\uPeel(\mathsf{IBLT}_0, \mathsf{IBLT}_1, i, j) = \begin{cases}
\mathsf{sum}_0[i][j], & \mathsf{cnt}_0{=}1 \land \mathsf{cnt}_1{=}0 \\
\mathsf{sum}_1[i][j], & \mathsf{cnt}_0{=}0 \land \mathsf{cnt}_1{=}1 \\
\mathsf{sum}_0[i][j], & \mathsf{cnt}_0{=}\mathsf{cnt}_1{=}1
                       \land \mathsf{sum}_0{=}\mathsf{sum}_1 \\
\bot, & \text{otherwise}
\end{cases}
\]
\end{definition}

Intuitively, $\uPeel$ distinguishes three cases: the peeled element is unique
to $P_0$ (case~1), unique to $P_1$ (case~2), or shared with equal sums (case~3).
Secure evaluation uses OT for sum retrieval and bOPRF for equality comparison.

\begin{theorem}[Union Peel Equivalence~\cite{PT26}]\label{thm:upeel}
For union-peelable IBLTs: $\Peel(\mathsf{IBLT}_\cup, i, j) =
\uPeel(\mathsf{IBLT}_0, \mathsf{IBLT}_1, i, j)$.
\end{theorem}

\begin{figure}[h]
\fbox{\begin{minipage}{0.95\columnwidth}
\smallskip
\noindent\textbf{Functionality $\cF_{\mathsf{UnionPeel}}$}

\noindent\textbf{Parameters:} $k, \ell, M$, hash functions $H = \{h_i\}_{i \in [k]}$.
$P_0$ holds $\mathsf{IBLT}_0 = \Encode(X_0)$,
$P_1$ holds $\mathsf{IBLT}_1 = \Encode(X_1)$.

\smallskip\noindent\textbf{Behavior:} $\cF_{\mathsf{UnionPeel}}$ behaves as follows.
\begin{enumerate}
\item Receive $\mathsf{IBLT}_0$ from $P_0$ and $\mathsf{IBLT}_1$ from $P_1$ (each
of which has been securely computed using ideal OT and OPRF). If either
$\mathsf{IBLT}$ is malformed return $\bot$ to both.
\item Compute the virtual union $\mathsf{IBLT}_\cup \leftarrow \mathsf{IBLT}_0
\oplus \mathsf{IBLT}_1$.
\item Run $\List(\mathsf{IBLT}_\cup)$. If peeling succeeds, send
$X_\cup = X_0 \cup X_1$ to both parties; otherwise send $\bot$.
\end{enumerate}
\smallskip
\end{minipage}}
\caption{The ideal functionality $\cF_{\mathsf{UnionPeel}}$ for secure union
peeling. Adapted from~\cite{PT26}.}
\label{fig:f-unionpeel}
\end{figure}

\begin{theorem}[Union Peel Security~\cite{PT26}]\label{thm:unionpeel-sec}
$\Pi_{\mathsf{UnionPeel}}$ UC-realizes $\cF_{\mathsf{UnionPeel}}$
(Fig.~\ref{fig:f-unionpeel}) in the $(\cF_{\mathsf{OPRF}},
\cF_{\mathsf{OT}})$-hybrid model against semi-honest adversaries, achieving
static PSU on $2^{20}$ elements in under 3~seconds.
\end{theorem}

\subsection{Cryptographic Primitives}

\textbf{ChaCha20~\cite{Ber08}.} We use ChaCha20 (256-bit key, 96-bit nonce) as
a secure PRG. Encryption: $C = M \oplus \ChaCha(K, N)$. The IBLT ciphertexts
are encrypted under per-epoch keys derived from the PRF chain.

\textbf{HMAC-SHA256~\cite{BCK96}.} We use HMAC-SHA256 as a secure PRF:
$F: \{0,1\}^\kappa \times \{0,1\}^* \to \{0,1\}^\kappa$. Key evolution uses a
fixed domain separator: $K_t \leftarrow F(K_{t-1}, \text{``evolve''})$.

\begin{lemma}[PRF Key One-Wayness]\label{lem:prf-ow}
If $F$ is a secure PRF, then for any PPT $\cA$ and fixed $c$:
$\Pr[\cA(F(K, c)) = K : K \sample \{0,1\}^\kappa] \leq \negl(\kappa)$.
\end{lemma}

\textbf{VOLE and OT extension.} Vector Oblivious Linear Evaluation (VOLE)
enables two parties to obtain additive shares of $\Delta \cdot \vec{x}$,
realizable from $\kappa$ base OTs with sublinear communication~\cite{RS21,RR22};
silent VOLE~\cite{RRT23} eliminates the base OTs entirely. The VOLE state is
reusable across all epochs, amortizing the one-time setup cost.

\textbf{Batched OPRF.} The KKRT16/RS21/RR22 framework~\cite{KKRT16,RS21,RR22}
constructs batched OPRF from VOLE and OKVS: the receiver encodes queries into
an OKVS $P$; VOLE shares $\Delta \cdot P$; OPRF outputs are defined as
$F_\Delta(x) = H(\Delta \cdot \mathsf{Decode}(P, x))$.

The ideal 1-out-of-2 OT functionality $\cF_{\mathsf{OT}}$ is defined in
Fig.~\ref{fig:f-ot}.

\begin{figure}[h]
\fbox{\begin{minipage}{0.95\columnwidth}
\smallskip
\noindent\textbf{Functionality $\cF_{\mathsf{OT}}$}

\noindent\textbf{Parameters:} Sender $S$ holds two messages
$(m_0, m_1) \in \{0,1\}^\kappa$; receiver $R$ holds a choice bit
$b \in \{0,1\}$.

\smallskip\noindent\textbf{Behavior:}
\begin{enumerate}
\item Receive $(m_0, m_1)$ from $S$ and $b$ from $R$. If either party
sends $\bot$, output $\bot$ to both and halt.
\item Send $m_b$ to $R$ and nothing to $S$.
\end{enumerate}
\smallskip
\end{minipage}}
\caption{The ideal 1-out-of-2 Oblivious Transfer functionality.}
\label{fig:f-ot}
\end{figure}

The ideal OPRF functionality $\cF_{\mathsf{OPRF}}$ is defined in
Fig.~\ref{fig:f-oprf}.

\begin{figure}[h]
\fbox{\begin{minipage}{0.95\columnwidth}
\smallskip
\noindent\textbf{Functionality $\cF_{\mathsf{OPRF}}$}

\noindent\textbf{Parameters:} Sender $S$ holds PRF key $k \in \{0,1\}^\kappa$;
receiver $R$ holds input $x \in \U$.

\smallskip\noindent\textbf{Behavior:}
\begin{enumerate}
\item Receive $k$ from $S$ and $x$ from $R$. If either party sends $\bot$,
output $\bot$ to both and halt.
\item Compute $y \leftarrow F(k, x)$. Send $y$ to $R$ and nothing to $S$.
\end{enumerate}
\smallskip
\end{minipage}}
\caption{The ideal Oblivious PRF functionality.}
\label{fig:f-oprf}
\end{figure}

The standard (static) PSU functionality $\cF_{\mathsf{PSU}}$ is defined
in Fig.~\ref{fig:f-psu}.

\begin{figure}[h]
\fbox{\begin{minipage}{0.95\columnwidth}
\smallskip
\noindent\textbf{Functionality $\cF_{\mathsf{PSU}}$}

\noindent\textbf{Parameters:} Two parties $P_0, P_1$ with element domain $\U$.

\smallskip\noindent\textbf{Behavior:}
\begin{enumerate}
\item Receive set $X$ from $P_0$ and $Y$ from $P_1$. If either is $\bot$,
output $\bot$ to both and halt.
\item Send $X \cup Y$ to both parties.
\end{enumerate}
\smallskip
\end{minipage}}
\caption{The ideal (static) Private Set Union functionality.}
\label{fig:f-psu}
\end{figure}

\subsection{UC Framework}

We use the Universal Composability framework~\cite{Can01} with adaptive
corruptions. A protocol $\Pi$ UC-realizes an ideal functionality $\cF$ if for
every PPT adversary $\cA$ there exists a PPT simulator $\cS$ such that no PPT
environment $\cZ$ can distinguish the real execution from the ideal execution.
We operate in the $(\cF_{\mathsf{OPRF}}, \cF_{\mathsf{OT}})$-hybrid model,
where the protocol invokes the ideal functionalities defined in
Figs.~\ref{fig:f-ot},~\ref{fig:f-oprf}. The standard (static) PSU
functionality $\cF_{\mathsf{PSU}}$ (Fig.~\ref{fig:f-psu}) and the
union-peeling functionality $\cF_{\mathsf{UnionPeel}}$
(Fig.~\ref{fig:f-unionpeel}) provide the building blocks. The ideal
functionality for our protocol, $\cF_{\mathsf{fpsu}}$, is specified in
Section~\ref{sec:protocol} (Fig.~\ref{fig:f-fpsu}).

% ============================================================
\section{Technical Overview}
% ============================================================

We illustrate F-PSU through a concrete epoch. Suppose $P_0$ holds a set $X$ and
$P_1$ holds $Y$, both having previously executed the protocol so that each
stores an encrypted IBLT encoding of their current set. At the start of a new
epoch, each party learns its delta: elements to add ($\Delta^+$) and delete
($\Delta^-$). The goal is to compute the updated union $X_\cup \leftarrow (X
\cup \Delta^+_X \setminus \Delta^-_X) \cup (Y \cup \Delta^+_Y \setminus
\Delta^-_Y)$ while revealing nothing beyond the union and the necessary leakage
identified above (item~4). F-PSU achieves this through three
coordinated mechanisms.

\smallskip\noindent\textbf{1. Delta IBLT encoding: $O(|\Delta|)$ local maintenance.}
By IBLT linearity (Algorithm~\ref{alg:encode}), each party maintains its IBLT
across epochs by applying only the delta:
$\mathsf{IBLT} \leftarrow \mathsf{IBLT} \ominus \Encode(\Delta^-) \oplus
\Encode(\Delta^+)$.
This costs $O(k|\Delta|)$ cell operations instead of $O(k|X|)$ for full
re-encoding. \emph{Deletions} are applied locally (cell-wise subtraction, zero
communication). \emph{Additions} are encoded into a separate compact IBLT
($\mathsf{IBLT}_{\Delta}$) but not yet merged---this deferred merge is the key
to avoiding cascade blow-up in the secure peel phase.

\smallskip\noindent\textbf{2. Split-path incremental peeling: $O(|\Delta|)$ secure computation.}
After each epoch, the main IBLT is fully peeled (no singletons remain). When
elements are deleted, the corresponding cells lose their contributions,
potentially creating new singletons. When elements are added, they would
introduce fresh singletons if merged directly. We avoid this interaction by
separating additions and deletions into two peeling paths:

\begin{itemize}
\item \textbf{Round~2a (Mini UnionPeel).} The compact addition IBLTs are
peeled via a fresh full UnionPeel~\cite{PT26} on only $m_{\Delta} =
c_k \cdot k \cdot |\Delta^+|$ cells. Since the addition IBLT has no prior
peeling state, every cell must be scanned, but the scan is deterministic:
zero cascade, zero variance. The recovered elements $V_{\Delta^+}$ are then
locally merged into the main IBLT (zero communication).

\item \textbf{Round~2b (Main UnionPeel).} For deletions, UnionPeel starts
\emph{only} from the $k|\Delta^-|$ cells modified by deletion (rather than
all $k\ell$ cells). The first iteration peels from these sites; recovered
elements are deleted from both IBLTs, and newly exposed singletons
(cascade) are added to the next query set. The \textbf{Cascade Bound Lemma}
(Lemma~\ref{lem:cascade}) proves that total examined cells is $\leq
2k|\Delta^-|$ with overwhelming probability.
\end{itemize}

Overall, per-epoch OT/OPRF cost is $O(k|\Delta^+| + k|\Delta^-|) =
O(k|\Delta|)$, a saving of $\approx n/|\Delta| \times$ over static
PSU~\cite{PT26} which scans all $k\ell$ cells.

\smallskip\noindent\textbf{3. Forward privacy via PRF chain: zero additional communication.}
Each party $P_b$ maintains a 256-bit key $K^b_t$, evolved via a PRF:
$K^b_t \leftarrow \PRF(K^b_{t-1}, \text{``evolve''})$. After each epoch, the
IBLT is encrypted under $K^b_t$ (ChaCha20), written to persistent storage, and
$K^b_{t-1}$ is immediately erased. By PRF one-wayness
(Lemma~\ref{lem:prf-ow}), an adversary compromising at epoch $\tau$ learns
$K^b_\tau$ and the current IBLT, but cannot recover any prior key
$K^b_{<\tau}$ or decrypt ciphertexts from $t < \tau$. The overhead is
negligible: one PRF evaluation plus one ChaCha20 encryption per epoch
($<1$~ms). This construction is the symmetric analogue of Signal's symmetric
ratchet~\cite{Signal}, applied to IBLT ciphertexts rather than message keys.

\smallskip\noindent\textbf{4. Leakage optimality: a matching lower bound.}\label{sec:barrier}

\begin{theorem}[Efficiency-Leakage Trade-off]\label{thm:lower}
Any correct, semi-honest secure, $1$-updatable PSU protocol necessarily reveals
$(|\Delta^+_0|, |\Delta^-_0|, |\Delta^+_1|, |\Delta^-_1|)$. Any protocol
hiding these counts must have $|\mathsf{trans}| = \Omega(N\sigma)$---defeating
updatability.
\end{theorem}

F-PSU leaks these four counts through observable delta IBLT sizes. The proof
uses a source-coding argument:
communication must vary with update size, otherwise an adversary could encode
arbitrary set information below the Shannon information-theoretic bound. Any
protocol hiding these counts requires $|\mathsf{trans}| = \Omega(N\sigma)$
---defeating updatability. F-PSU leaks exactly these four counts and
nothing more; by Theorem~\ref{thm:lower}, it is \textbf{strictly
leakage-optimal}. The full proof appears in Appendix~\ref{sec:lower-proof}.

\smallskip\noindent\textbf{Security.}
Per-epoch security inherits from $\Pi_{\mathsf{UnionPeel}}$~\cite{PT26}.
Cross-epoch forward privacy is proved via a three-hybrid argument
(Section~\ref{sec:security}): (i)~replace PRF-derived keys with random keys
($T \cdot \mathsf{Adv}^{\mathsf{PRF}}$ reduction), (ii)~replace IBLT
ciphertexts with encryptions of zero ($T \cdot \mathsf{Adv}^{\mathsf{PRG}}$
reduction), (iii)~replace real inputs with simulated inputs derived from the
union output. The modular proof structure enables independent upgrades of the
per-epoch PSU and the cross-epoch key management.

% ============================================================
\section{System Model and Threat Model}
% ============================================================

\textbf{System model.} Two parties $P_0, P_1$ hold private sets $X_0, Y_0
\subseteq \{0,1\}^\sigma$. The protocol runs over $T = \poly(\kappa)$ epochs.
At epoch $t$, each party provides a delta: elements to add ($\Delta^+$) and
delete ($\Delta^-$), disjoint from each other and from the current set. After
epoch $t$, both parties learn the updated union $X_{\cup,t}$.

\textbf{Updatable efficiency.} A protocol is \emph{$1$-updatable} if per-epoch
communication $|\mathsf{trans}_t| = O(|\Delta_t| \cdot \poly(\sigma, \kappa))$,
scaling with delta size rather than total set size.

\textbf{Threat model.} Semi-honest adversary with adaptive corruption of one
party at epoch $\tau$ (corrupting both is trivial). The adversary observes all
transcripts. Upon corruption, it learns the corrupted party's current state
(set $X_\tau$, key $K_\tau$, IBLT, VOLE/OKVS state). Forward privacy requires
learning \emph{nothing} about sets from epochs $t < \tau$ beyond what union
outputs already reveal. Post-compromise security is not required in the main
protocol; a self-healing extension appears in Appendix~\ref{sec:selfheal}.

The formal ideal functionality $\cF_{\mathsf{updPSO}}$
(Fig.~\ref{fig:f-fpsu}) and UC definitions appear in
Sections~\ref{sec:protocol} and~\ref{sec:security}.
$\cF_{\mathsf{updPSO}}$ receives deltas, computes the updated union and
intersection, and returns the output corresponding to a mode parameter
($\mathsf{PSU}, \mathsf{PSI}, \mathsf{PSU}\text{-}\mathsf{CA},
\mathsf{PSI}\text{-}\mathsf{CA}$). Theorem~\ref{thm:lower} proves the
leaked delta counts are unavoidable.

% ============================================================
\section{The F-PSU Protocol}
\label{sec:protocol}

\subsection{Ideal Functionality}

\begin{figure}[h]
\fbox{\begin{minipage}{0.95\columnwidth}
\smallskip
\noindent\textbf{Functionality $\cF_{\mathsf{updPSO}}$}

\noindent\textbf{Parameters:} Two parties $P_0, P_1$, security parameter
$\kappa$, element domain $\U$, max set size $N$, mode
$\mathsf{mode} \in \{\textsf{PSU}, \textsf{PSI}, \textsf{PSU}\text{-}\textsf{CA},
\textsf{PSI}\text{-}\textsf{CA}\}$. For each $P_b$, the
functionality maintains state $(X_b, K_b)$.

\smallskip\noindent\textbf{Behavior:} $\cF_{\mathsf{updPSO}}$ behaves as follows.
\begin{enumerate}
\item Upon receiving $(\textsf{Update}, \Delta^+_b, \Delta^-_b)$ from $P_b$:
update $X_b \leftarrow (X_b \cup \Delta^+_b) \setminus \Delta^-_b$;
send $(|X_b|, |\Delta^+_b|, |\Delta^-_b|)$ to the simulator $\cS$.
\item Once both parties have submitted their updates, compute $X_\cup = X_0 \cup
X_1$ and $X_\cap = X_0 \cap X_1$. If $t=1$, sample $K_b \sample
\{0,1\}^\kappa$ for each $b$; otherwise $K_b \leftarrow \PRF(K_b)$ and
immediately erase the old $K_b$. Send the output according to
$\mathsf{mode}$:
  \begin{itemize}
  \item $\mathsf{mode} = \textsf{PSU}$: send $X_\cup$ to both parties.
  \item $\mathsf{mode} = \textsf{PSI}$: send $X_\cap$ to both parties.
  \item $\mathsf{mode} = \textsf{PSU}\text{-}\textsf{CA}$: send $|X_\cup|$
  to both parties.
  \item $\mathsf{mode} = \textsf{PSI}\text{-}\textsf{CA}$: send $|X_\cap|$
  to both parties.
  \end{itemize}
\item Upon adaptive corruption of $P_b$ at epoch $\tau$: send
$(X_b^{(\tau)}, K_b^{(\tau)})$ to $\cS$. \textbf{Not} sent:
$\{(X_b^{(t)}, K_b^{(t)})\}_{t < \tau}$.
\end{enumerate}
\smallskip\noindent\textbf{Leakage:} $\cL = (|X_0|, |X_1|, |\Delta^+_0|,
|\Delta^-_0|, |\Delta^+_1|, |\Delta^-_1|)$ to the simulator.
\smallskip
\end{minipage}}
\caption{The ideal functionality $\cF_{\mathsf{updPSO}}$ for forward-private
updatable Private Set Operations (PSU, PSI, and their cardinalities).}
\label{fig:f-fpsu}
\end{figure}

% ============================================================

\subsection{Extension: PSI and Cardinality from Pure Cells}

The $\uPeel$ function (Definition~\ref{def:upeel}) naturally reveals whether a
peeled element belongs to the intersection. When $\uPeel$ peels a cell
$(i,j)$, it distinguishes three cases: (1)~$\mathsf{cnt}_0[i][j]=1,
\mathsf{cnt}_1[i][j]=0$---the element is unique to $P_0$;
(2)~$\mathsf{cnt}_0[i][j]=0, \mathsf{cnt}_1[i][j]=1$---unique to $P_1$;
(3)~$\mathsf{cnt}_0[i][j]=1, \mathsf{cnt}_1[i][j]=1, \mathsf{sum}_0=
\mathsf{sum}_1$---the element belongs to \textbf{both} parties, i.e., the
intersection $X_0 \cap X_1$. Case~(3) occurs because each party independently
encodes its elements into its own IBLT; when the same element $x$ is encoded
by both, it lands in the same cell $(i,j)$ with identical sum $x$, and
$\uPeel$ detects the match.

Consequently, a single execution of the UnionPeel sub-protocol simultaneously
recovers the union (all peeled elements) and the intersection (Case~3
elements). This yields four output modes at zero additional communication:
\begin{itemize}
\item \textbf{PSU:} output all peeled elements ($X_0 \cup X_1$).
\item \textbf{PSI:} output only Case~3 elements ($X_0 \cap X_1$).
\item \textbf{PSU-CA:} output the count $|X_0 \cup X_1|$.
\item \textbf{PSI-CA:} output the count $|X_0 \cap X_1|$.
\end{itemize}

The mode is selected via a single bit in the protocol initialization and
requires no change to the protocol messages, the OT/OPRF interaction, the
IBLT encoding, or the peeling logic. The only modification is that the
peeling output carries a 2-bit tag per element identifying the case, and the
parties filter output accordingly.

Security inherits directly from Theorem~\ref{thm:main}: the simulator for
$\Pi_{\mathsf{fpsu}}$ already extracts the corrupted party's effective input
and submits it to the ideal functionality. The same simulator, paired with
$\cF_{\mathsf{updPSO}}$ (which computes the output for the chosen mode),
produces an indistinguishable view. The per-mode leakage is identical:
$X_\cup$ reveals $Y \setminus X$ in all modes; $X_\cap$ additionally reveals
element equality between the two sets; cardinalities reveal only the sizes.
The adaptive security argument of Section~\ref{sec:security} applies
unchanged since the protocol transcript is independent of the output mode.

\subsection{Initialization}

At setup, $P_0$ and $P_1$ (denoted Client $C$ and Server $S$ in pseudocode)
perform one-time VOLE setup (RR22~\cite{RR22}) reused across all epochs.
Each party initializes an empty IBLT of capacity $(m, k)$ and samples an
independent 256-bit PRF key.

\subsection{Epoch Protocol}

Each epoch comprises three sequential rounds, with Round~2 split into two
phases (2a for additions, 2b for deletions). Pseudocode follows the
complexity analysis below.

\noindent\textbf{Round~1 (Local IBLT update).}
Each party decrypts its stored IBLT. Deletions are applied directly via
cell-wise subtraction ($\mathsf{IBLT} \leftarrow \mathsf{IBLT} \ominus
\Encode(\Delta^-)$), touching exactly $k|\Delta^-|$ cells with zero
communication. The coordinates of modified cells
are recorded: $Q_{\mathsf{del}} = \{(i, h_i(x)) : x \in \Delta^-, i \in [k]\}$.
Additions are encoded into a separate compact IBLT but \emph{not} merged yet.

\noindent\textbf{Round~2a (Mini UnionPeel for additions).}
$\Pi_{\mathsf{UnionPeel}}$ on the compact addition IBLTs scans all $m_{\Delta}
= c_k \cdot k \cdot |\Delta^+|$ cells. With no prior peeling state, this is
deterministic with zero cascade, recovering $V_{\Delta^+}$ with probability
$\geq 1 - 2^{-\lambda}$. After recovery, each party locally merges
$\Encode(\Delta^+)$ into the main IBLT.

\noindent\textbf{Round~2b (Main UnionPeel for deletions).}
$\Pi_{\mathsf{UnionPeel}}$ on the main IBLTs starts from $Q^{(0)} =
Q_{\mathsf{del}}^C \cup Q_{\mathsf{del}}^S$ ($k|\Delta^-|$ cells). The
procedure iteratively evaluates $\uPeel$ via OT and bOPRF, deletes peeled
elements, and propagates the cascade. By Lemma~\ref{lem:cascade}, $\leq
2k|\Delta^-|$ cells are examined. $P_1$ sends $V_{\mathsf{main}}$ to $P_0$;
both output $X_{\cup,t} = V_{\Delta^+} \cup V_{\mathsf{main}}$.

\noindent\textbf{Round~3 (Key evolution).}
$K^b_t \leftarrow \PRF(K^b_{t-1})$, encrypt IBLT under $K^b_t$ (ChaCha20),
write ciphertext, \textbf{securely erase} $K^b_{t-1}$.

\noindent\textbf{Why split-path?} Without separation, addition cells would mix
into the deletion cascade, growing it proportionally to $d^+ + d^-$. The mini
peel gives additions a dedicated deterministic channel; deletions are handled
incrementally. This separation is the algorithmic core reducing OT/OPRF cost
from $O(n)$ to $O(|\Delta|)$.

\begin{figure}[h]
\fbox{\begin{minipage}{0.95\columnwidth}
\smallskip
\noindent\textbf{Protocol $\Pi_{\mathsf{fpsu}}$ (epoch $t$)}

\smallskip\noindent\textit{Public parameters:} IBLT parameters $\prm = (M, k,
\ell, H)$, PRF $F$, stream cipher $\ChaCha$, expansion factor $c_k$.

\noindent\textit{Party state:} $P_b$ holds encrypted $\mathsf{IBLT}_b$,
current set $X_b$, key $K^b_{t-1}$.

\smallskip\noindent\textbf{Each epoch:}
\begin{enumerate}
\item \textbf{Local IBLT update (Round~1).} Decrypt $\mathsf{IBLT}_b$;
$\mathsf{IBLT}_b \leftarrow \mathsf{IBLT}_b \ominus \Encode(\Delta^-_b)$;
$\mathsf{IBLT}^+_b \leftarrow \Encode(\Delta^+_b)$ (compact).
Record $Q_{\mathsf{del}} = \{(i, h_i(x)) : x \in \Delta^-_b, i \in [k]\}$.

\item \textbf{Mini UnionPeel for additions (Round~2a).} Run
$\Pi_{\mathsf{UnionPeel}}$ on $\mathsf{IBLT}^+_C, \mathsf{IBLT}^+_S$
(scanning $m_{\Delta} = c_k \cdot k \cdot \max(|\Delta^+_C|,
|\Delta^+_S|)$ cells). Recover $V_{\Delta^+}$, merge:
$\mathsf{IBLT}_b \leftarrow \mathsf{IBLT}_b \oplus \Encode(V_{\Delta^+})$.

\item \textbf{Main UnionPeel for deletions (Round~2b).} Run
$\Pi_{\mathsf{UnionPeel}}$ on $\mathsf{IBLT}_C, \mathsf{IBLT}_S$
starting from $Q_{\mathsf{del}}^C \cup Q_{\mathsf{del}}^S$. Recover
$V_{\mathsf{main}}$, output $X_{\cup,t} = V_{\Delta^+} \cup
V_{\mathsf{main}}$. If $\mathsf{mode} = \mathsf{PSI}$, output only Case~3
elements.

\item \textbf{Key evolution (Round~3).} $K^b_t \leftarrow \PRF(K^b_{t-1},
\text{``evolve''})$; $\mathsf{CT}_b \leftarrow \ChaCha.\Enc(K^b_t,
\mathsf{IBLT}_b)$; write $\mathsf{CT}_b$; \textbf{erase} $K^b_{t-1}$.
\end{enumerate}
\smallskip
\end{minipage}}
\caption{The F-PSU epoch protocol $\Pi_{\mathsf{fpsu}}$.}
\label{fig:protocol}
\end{figure}

\subsection{Complexity}

\textbf{Unbalanced settings.} F-PSU operates identically in balanced and
unbalanced scenarios because IBLT encoding complexity depends only on the
\emph{local} set size, not the peer's set size. When $P_A$ holds $2^{10}$
elements and $P_B$ holds $2^{20}$, each party encodes its own IBLT with cost
proportional to $|\text{local set}|$, independent of the other. The delta
IBLT mechanism is similarly local: each party applies $\Delta$ only to its
own IBLT. The UnionPeel sub-protocol treats both IBLTs symmetrically,
peeling from the virtual union regardless of size disparity. No protocol
modification is required. The asymptotic complexity is $O(n\sigma)$ for the
union output and $O(d(\sigma + \kappa))$ for the per-epoch update, where
$n = |X| + |Y|$ and $d = |\Delta|$.

Let $d = |\Delta| = d^+ + d^-$, $n = |X| + |Y|$.
\begin{itemize}
\item Local update: $O(k d)$ cell ops, zero communication.
\item Mini UnionPeel: $O(c_k \cdot k \cdot d^+ \cdot (\sigma + \kappa))$ bits,
deterministic, zero cascade.
\item Main UnionPeel: $O(k d^- \cdot (\sigma + \kappa))$ bits, cascade $\leq
2k d^-$ cells.
\item Output: $n\sigma$ bits (inherent to any PSU).
\item Forward privacy overhead: one PRF + one ChaCha20 per epoch ($<1$~ms).
\end{itemize}
Total: $O(d(\sigma + \kappa) + n\sigma)$. For $d/n \approx 0.1\%$, the non-output
component is $\approx 1000\times$ smaller than static PSU~\cite{PT26}.

\begin{table}[t]
\centering
\footnotesize
\caption{Asymptotic comparison of PSU protocols ($d = |\Delta|$, $n = |X|+|Y|$).}
\label{tab:complexity}
\footnotesize
\resizebox{\columnwidth}{!}{%
\begin{tabular}{lccc}
\toprule
\textbf{Protocol} & \textbf{Communication} & \textbf{Computation} & \textbf{Bal./Unbal.} \\
\midrule
KRT19~\cite{KRT19} & $O(n\sigma\kappa)$ & $O(n\log n)$ OT & Balanced \\
PT26~\cite{PT26} & $O(n(\sigma+\kappa))$ & $O(n)$ OT+hash & Balanced \\
ePSU~\cite{Tu+25} & $O(n\sigma)$ & $O(n)$ symm. & Both \\
\textbf{F-PSU} & $\mathbf{O(d(\sigma+\kappa)+n\sigma)}$ & $\mathbf{O(d)}$ OT$+\mathbf{O(n)}$ ChCha & \textbf{Both} \\
\bottomrule
\end{tabular}%
}
\end{table}

\begin{lemma}[Cascade Bound]\label{lem:cascade}
With IBLT parameters ensuring 2-core emptiness ($\Pr \geq 1 - 2^{-\lambda}$),
peeling from $k d^-$ deletion-modified cells examines at most $(1 + C_0) \cdot
k d^-$ total cells, where $C_0 = O(k^2 d^- / m)$. For $d^- \ll m/k$, the
cascade is $\leq 2k d^-$ with overwhelming probability.
\end{lemma}

\subsection{Communication Reduction}

\begin{figure}[t]
\centering
\includegraphics[width=\columnwidth]{fig_communication.pdf}
\caption{Per-epoch IBLT data exchange. Left: static vs.\ F-PSU (log scale,
$|\Delta|/n = 0.1\%$). Right: communication reduction factor across set sizes.}
\label{fig:comm}
\end{figure}

Table~\ref{tab:main-perf} compares per-epoch IBLT data exchange against static
re-encoding. At $0.1\%$ update ratio, communication is reduced by over
$643\times$. The reduction is exactly $n/|\Delta|$---the ratio of total set
size to delta size---because F-PSU exchanges only the cells touched by
$\Delta$, while static protocols exchange all cells. At $10\%$ update rate
the reduction drops to $6.3\times$, matching the expected $1/0.1 = 10\times$
ratio with IBLT expansion overhead. OT/OPRF communication scales
proportionally with cell count, so the same reduction factor applies to the
secure computation phase.



\subsection{Incremental Peeling and Parameter Growth}

For $n = 2^{16}$, $|\Delta| = 2^{12}$: incremental peeling examines 32,768
cells vs.\ 197,000 for full scan ($6\times$ reduction). Per-epoch protocol:
$20.8$~ms ($n = 2^{14}$, $|\Delta| = 2^{12}$), $97.8$~ms ($n = 2^{16}$,
$|\Delta| = 2^{14}$). Forward privacy overhead: $<1$~ms.

\textbf{Amortized parameter growth.} 200 epochs, $n_0 = 10^4$, 100 elements
per epoch: only \textbf{1} full re-encode triggered (epoch~51).

\begin{figure}[t]
\centering
\includegraphics[width=\columnwidth]{fig_amortized.pdf}
\caption{Amortized IBLT cell operations over 200 epochs, $|\Delta|/n = 0.1\%$.
F-PSU achieves $114\times$ reduction across all set sizes.}
\label{fig:amortized}
\end{figure}

% ============================================================
\section{Security Analysis}
\label{sec:security}
% ============================================================

We prove that $\Pi_{\mathsf{fpsu}}$ UC-realizes $\cF_{\mathsf{updPSO}}$
(Fig.~\ref{fig:f-fpsu}) in the $(\cF_{\mathsf{OPRF}}, \cF_{\mathsf{OT}})$-hybrid
model against PPT semi-honest adaptive adversaries, for all four output modes.

\subsection{UC Security}

\begin{theorem}[Main Security]\label{thm:main}
Assume ChaCha20 is a secure PRG and HMAC-SHA256 is a secure PRF. Then
$\Pi_{\mathsf{fpsu}}$ UC-realizes $\cF_{\mathsf{updPSO}}$ in the
$(\cF_{\mathsf{OPRF}}, \cF_{\mathsf{OT}})$-hybrid model against PPT semi-honest
adaptive adversaries.
\end{theorem}

\textbf{Proof structure.} The proof decomposes modularly into per-epoch security
(within an epoch, inheriting from $\Pi_{\mathsf{UnionPeel}}$~\cite{PT26}) and
cross-epoch forward privacy (via the PRF chain). The full proof, including
simulator construction and three-hybrid argument, appears in
Appendix~\ref{sec:full-proofs}.

\subsection{Game-Based Forward Privacy (FP-IND)}

We supplement the UC definition with a standalone game-based notion that
directly captures the PRF-chain guarantee.

\begin{definition}[FP-IND Game]\label{def:fp-ind}
\begin{enumerate}
\item Challenger samples $b \sample \{0,1\}$, $K_0 \sample \{0,1\}^\kappa$.
\item For $t = 1$ to $T$: $K_t \leftarrow F(K_{t-1}, c)$.
\item Adversary (observing all transcripts) selects $\tau \in [T]$, receives
$K_\tau$.
\item If $b = 0$: receives real $\{K_t\}_{t < \tau}$. If $b = 1$: receives
independent $\{R_t \sample \{0,1\}^\kappa\}_{t < \tau}$.
\item Adversary outputs $b'$, wins if $b' = b$.
\end{enumerate}
$\mathsf{Adv}^{\mathsf{FP\text{-}IND}}(\kappa) = |\Pr[b' = b] - \tfrac{1}{2}|
\leq \negl(\kappa)$.
\end{definition}

\begin{corollary}[Forward Privacy]\label{cor:fp}
If $\Pi_{\mathsf{fpsu}}$ UC-realizes $\cF_{\mathsf{updPSO}}$ with adaptive
corruptions, then $\mathsf{Adv}^{\mathsf{FP\text{-}IND}}(\kappa) \leq T \cdot
\mathsf{Adv}^{\mathsf{PRF}}(\kappa)$.
\end{corollary}

\subsection{Adaptive Input Selection}

Liu et al.~\cite{LMR+26} recently showed that several published UPSI protocols
are broken by \emph{adaptive input selection}: an adversary who chooses deltas
based on prior transcripts can extract information about the honest party's
historical inputs from the protocol's stored state (e.g., OKVS encodings). We
show that F-PSU requires no modification to resist this class of attacks.

\smallskip\noindent\textbf{Why UPSI is vulnerable.}
In LMR+26's setting, the server stores an OKVS encoding of its set across
epochs. An adaptive adversary can insert elements strategically chosen to
interact with this encoding, learning membership information beyond what the
PSI output reveals. The root cause is that the server's \emph{state} depends
deterministically on its historical inputs, and adaptive queries can probe
this state.

\smallskip\noindent\textbf{Why F-PSU is immune.}
F-PSU's state has two layers of protection:
\begin{enumerate}
\item \textbf{Semantic security of encryption.} The stored IBLT is encrypted
under ChaCha20 with a per-epoch key $K_t$. An adaptive adversary probing the
ciphertext learns nothing about the underlying plaintext (semantic security of
ChaCha20 as a PRG).
\item \textbf{Forward privacy of the key chain.} The key $K_t$ is derived via a
PRF chain: $K_t \leftarrow \PRF(K_{t-1})$, with $K_{t-1}$ immediately erased.
By the FP-IND guarantee (Corollary~\ref{cor:fp}), compromising at epoch $\tau$
reveals $K_\tau$ but nothing about $\{K_t\}_{t < \tau}$. The adversary cannot
decrypt ciphertexts from prior epochs, even with adaptive oracle access.
\end{enumerate}

Consequently, the adversary's view at each epoch consists of the union output
$X_{\cup,t}$ and the size leakage $\cL_t$, exactly as specified by the ideal
functionality $\cF_{\mathsf{fpsu}}$ (Fig.~\ref{fig:f-fpsu}). Adaptive choice
of deltas changes \emph{which} elements are added or deleted, but does not
change \emph{what the protocol leaks}: the updated union and the delta sizes.
No additional game-based definition is required; the UC simulator in
Theorem~\ref{thm:main} already handles adaptive input selection by extracting
the corrupted party's effective input and forwarding it to $\cF_{\mathsf{fpsu}}$.

\smallskip\noindent\textbf{Structural contrast with PSI.}
In PSI, the output $X_{\cap,t}$ is a \emph{subset} of both inputs: an element
appears in the output only if both parties hold it. Adaptive addition of $x$
reveals $x \in Y_t$ with certainty (if $x$ appears in the new intersection).
In PSU, the output $X_{\cup,t}$ is a \emph{superset}: any element added by the
adversary appears in the union unconditionally. The adversary learns
$Y_t \setminus X_{0,t}$ (inherent to PSU), but adaptive probing yields no
additional information---the encrypted state and forward-private keys ensure
that historical data remains inaccessible regardless of input strategy.

% ============================================================
\section{Implementation and Evaluation}
\label{sec:eval}
% ============================================================

We implement F-PSU in \textbf{C++} ($\approx$2,000 lines, six modules) on an
Apple M5 (macOS, single-threaded) using libOTe~\cite{libOTe} v2.1. IBLT
parameters: $k = 4$, expansion $e = 3.0$, $\sigma = 64$ bits. Network
conditions beyond LAN are simulated in-app with configurable RTT and bandwidth
caps. Code:~\cite{code}.

\subsection{Main Performance Comparison}

\begin{table*}[t]
\centering
\footnotesize
\caption{Performance comparison with prior PSU protocols. Baseline data
(ZCL+23, Jia+24, Tu+25, PT26) from~\cite{PT26} (AMD EPYC 74F3,
single-threaded). F-PSU on Apple M5 (single-threaded).
F-PSU epoch~0 = full IBLT re-encode; epoch~$>$0 = delta IBLT at
$|\Delta|/n \approx 6.25\%$. Bold = best; ``---'' = not reported.
F-PSU communication = IBLT data exchange (dominant per-epoch cost;
OT/OPRF setup amortized). Cross-platform runtime comparison is
approximate; the primary comparison is F-PSU epoch~0 vs.\ epoch~$>$0
on identical hardware.}
\label{tab:main-perf}
\begin{tabular}{lcccccc}
\toprule
$\mathbf{n}$ & \textbf{Protocol} & \textbf{Comm. (MB)} & \multicolumn{4}{c}{\textbf{Runtime (sec)}} \\
             &                   &                    & LAN & WAN 200M & WAN 50M & WAN 5M \\
\midrule
\multirow{6}{*}{$2^{14}$}
& ZCL+23 (static)      & 2.73  & 24.15 & 40.99 & 42.15 & 125.7 \\
& Jia+24 (static)      & 32.9  & 1.26  & 13.97 & 17.53 & 285.9 \\
& Tu+25 (static)       & 4.76  & 2.04  & 5.63  & 5.93  & 43.19 \\
& PT26 (static)        & 11.9  & 0.08  & 2.42  & 3.85  & 79.2 \\
& F-PSU epoch~0        & 0.58$^\ddagger$  & 0.033 & ---   & 0.54  & 4.15 \\
& F-PSU epoch~$>$0     & 0.080$^\ddagger$ & \textbf{0.020} & 0.084 & \textbf{0.16} & \textbf{0.40} \\
\midrule
\multirow{6}{*}{$2^{16}$}
& ZCL+23 (static)      & 11.1  & 77.16 & 139.8 & 166.6 & 522.3 \\
& Jia+24 (static)      & 67.8  & 2.06  & 44.10 & 82.74 & 1291 \\
& Tu+25 (static)       & 9.18  & 20.16 & 78.54 & 83.91 & 452.2 \\
& PT26 (static)        & 10.7  & 0.14  & 4.22  & 7.60  & 161.2 \\
& F-PSU epoch~0        & 2.30$^\ddagger$  & 0.078 & 0.50  & 1.54  & 13.58 \\
& F-PSU epoch~$>$0     & 0.32$^\ddagger$  & \textbf{0.097} & 0.17 & \textbf{0.31} & \textbf{1.23} \\
\midrule
\multirow{6}{*}{$2^{18}$}
& ZCL+23 (static)      & 47.2  & 260.8 & 427.6 & 392.7 & 1567 \\
& Jia+24 (static)      & 166   & 5.43  & 131.1 & 276.1 & 4436 \\
& Tu+25 (static)       & 12.7  & 130.3 & 529.1 & 464.7 & 3734 \\
& PT26 (static)        & 10.2  & 0.21  & 8.49  & 12.33 & 248.3 \\
& F-PSU epoch~0        & 9.22$^\ddagger$  & 0.52  & 1.47  & 4.86  & 43.4 \\
& F-PSU epoch~$>$0     & 1.27$^\ddagger$  & \textbf{0.56} & 0.62 & \textbf{1.00} & \textbf{4.66} \\
\midrule
\multirow{6}{*}{$2^{20}$}
& ZCL+23 (static)      & 191   & 852.9 & ---   & ---   & --- \\
& Jia+24 (static)      & 373   & 22.57 & 517.3 & 1143  & 19313$^\dagger$ \\
& Tu+25 (static)       & 25.0  & 1073  & ---   & ---   & --- \\
& PT26 (static)        & 10.2  & 2.95  & 30.19 & 20.14 & --- \\
& F-PSU epoch~0        & 36.9$^\ddagger$  & 2.13  & 4.96  & 18.97 & --- \\
& F-PSU epoch~$>$0     & 5.09$^\ddagger$  & \textbf{2.28} & \textbf{2.49} & \textbf{3.40} & --- \\
\bottomrule
\end{tabular}
\end{table*}

\smallskip\noindent
At $0.1\%$ update ratio ($\approx$1,000 elements changed per epoch at
$n=2^{20}$), F-PSU epoch~$>$0 communication drops to $57.3$~KB
($643\times$ reduction vs.\ epoch~0), with IBLT computation $<5$~ms.
F-PSU is the only protocol supporting updatable execution with forward
privacy. The per-epoch breakdown across network profiles appears in the
F-PSU epoch~$>$0 rows above.

IBLT peeling dominates epoch time ($>85\%$). Forward privacy overhead
(one PRF + one ChaCha20 per epoch) is $<1$~ms in all configurations.



\subsection{Summary of Experimental Findings}

\begin{enumerate}
\item \textbf{Cascade bound tightness.} Cascade $\approx 1.95 \times
k|\Delta|$, within the $2 \times k|\Delta|$ bound. Ratio is stable across all
parameters.
\item \textbf{Communication reduction.} $643\times$ at $0.1\%$ updates,
decreasing gracefully to $6.3\times$ at $10\%$, matching $n/|\Delta|$ scaling.
\item \textbf{Network-adaptive benefit.} $2\times$ to $11\times$ speedup across
bandwidth-constrained profiles. On LAN, no wall-clock benefit---as expected.
\item \textbf{Negligible forward privacy overhead.} One PRF + one ChaCha20:
$<1$~ms in all configurations.
\end{enumerate}

\subsection{Feature Comparison}

Table~\ref{tab:compare} situates F-PSU in the landscape. No prior work
simultaneously achieves updatability, forward privacy, and unified PSO
(PSU + PSI + cardinalities) in a single protocol execution.

\begin{table*}[t]
\centering
\footnotesize
\caption{Feature comparison with prior PSU and updatable protocols.
$\dagger$: enhanced security (no within-execution leakage);
$\ddagger$: PSI; $\mathsection$: unified PSO (PSU+PSI+CA in one execution).}
\label{tab:compare}
\footnotesize
\begin{tabular}{lccccccc}
\toprule
\textbf{Protocol} & \textbf{Type} & \textbf{Upd.} & \textbf{FW} & \textbf{UC} & \textbf{Adp.} & \textbf{Unified} & \textbf{Cost/Epoch} \\
\midrule
KRT19~\cite{KRT19} & PSU (2p) & --- & --- & \checkmark & --- & --- & $O(n)$ \\
ZCL+23~\cite{ZCL+23} & PSU (2p) & --- & --- & \checkmark & --- & --- & $O(n)$ \\
Jia+24~\cite{Jia+24} & PSU (2p) & --- & ---$^\dagger$ & --- & --- & --- & $O(n\log n)$ \\
Tu+25~\cite{Tu+25} & PSU (2p) & --- & ---$^\dagger$ & \checkmark & --- & --- & $O(n)$ \\
PT26~\cite{PT26} & PSU (2p) & --- & --- & \checkmark & --- & --- & $O(n)$ \\
FUPSI~\cite{WZC+24} & PSI$^\ddagger$ (2p) & \checkmark & \checkmark & --- & --- & --- & $O(\log n)$ PIR \\
LMR+26~\cite{LMR+26} & PSI$^\ddagger$ (2p) & \checkmark & --- & \checkmark & \checkmark & --- & $O(|\Delta|)$ \\
\midrule
\textbf{F-PSU (this work)} & \textbf{PSU (2p)} & \textbf{\checkmark} & \textbf{\checkmark} & \textbf{\checkmark} & \textbf{\checkmark} & $\mathsection$ & $\mathbf{O(|\Delta|)}$ \\
\bottomrule
\end{tabular}
\end{table*}

\subsection{Deployment Perspective}

To ground these numbers: in threat-intelligence sharing with $n = 2^{20}
\approx 10^6$ IoCs and $0.1\%$ daily churn ($\approx$1,000 new/retired
indicators per day), F-PSU exchanges $\approx$57~KB of IBLT data per daily
epoch vs.\ $\approx$36~MB for static. Over a year: $\approx$13~GB saved per
party pair. In federated learning with 100~clients and $n = 10^6$ features,
daily PSU-based feature alignment would save several terabytes per year. The
per-epoch computational overhead ($<5$~ms of IBLT operations at this update
ratio) is negligible compared to network latency.

\textbf{Implementation scope.} Six C++ modules: IBLT data structure (319
lines), protocol orchestration (226 lines), TCP socket communication (271
lines), end-to-end network benchmark ($\approx$400 lines), OT/OPRF-integrated
benchmark ($\approx$330 lines), primitive benchmarks ($\approx$500 lines). All
IBLT operations achieve $100\%$ success across $2^{10}$ to $2^{20}$. Results
are reproducible on modern macOS without specialized hardware.

% ============================================================
\section{Related Work}
\label{sec:related}
% ============================================================

\textbf{Static PSU.} Early PSU constructions relied on additively homomorphic
encryption; KRT19~\cite{KRT19} pioneered scalable symmetric-key PSU via
reversed private membership test. JLZ20~\cite{JLZ20} (USENIX'22) proposed shuffle-based
PSU with symmetric-key operations, but $O(n \log n)$ communication.
ZCL+23~\cite{ZCL+23} (USENIX'23) achieved the first linear-complexity PSU via
multi-query RPMT. Jia+24~\cite{Jia+24} (USENIX'24) identified unnecessary OT
leakage in existing scalable PSU and built an enhanced-security PSU from
symmetric-key operations, though with superlinear complexity. Tu+25~\cite{Tu+25}
(USENIX'25) resolved the open problems from Jia+24: the first linear enhanced
PSU (balanced, 2.2--8.8$\times$ communication reduction) and the first unbalanced
enhanced PSU (sublinear communication, 38.1$\times$ reduction at 1~Mbps).
LBL25~\cite{LBL25} further eliminated residual privacy leakage from
Tu+25's hashing-to-bin step, achieving the first leakage-free unbalanced ePSU.
Dong+25~\cite{Dong+25} (USENIX'25) extended PSU to the multi-party setting with
linear complexity. PT26~\cite{PT26} (Eurocrypt'26) closed the PSI-PSU
performance gap via IBLT union peelability, achieving PSU on $2^{20}$ elements
in under 3 seconds. Pu et al.~\cite{PuGT26} (Eurocrypt'26) extended this to
malicious security. All prior PSU protocols---including the IBLT-based
constructions---are \emph{static}: they re-encode the entire set each epoch,
costing $O(n)$ per computation. \textbf{F-PSU is the first updatable PSU.}

\textbf{Updatable protocols} (PSI only). BMX22~\cite{BMX22} initiated updatable
PSI from OKVS. BMS+24~\cite{BMS+24} extended to support deletion with
worst-case bounds. LMR+26~\cite{LMR+26} proposed symmetric-key updatable PSI,
identifying adaptive input selection attacks that break prior UPSI proofs.
LTS+26~\cite{LTS+26} combined asymmetric PSI and PSU for low-communication
updatable PSI. FUPSI~\cite{WZC+24} (TIFS'24) achieved forward privacy in
updatable PSI via keyword PIR, with $O(\log n)$ communication per update.
\textbf{No prior work supports updatable PSU}, let alone with forward privacy.

\textbf{Forward-secure and privacy-preserving protocols.} Signal's double
ratchet~\cite{Signal} provides forward secrecy via Diffie-Hellman key exchange.
Bost~\cite{Bos16} achieved forward privacy for searchable symmetric encryption.
Our PRF-based key chain is the symmetric analogue of Signal's symmetric ratchet,
with the key distinction that we operate on IBLT ciphertexts rather than session
keys. The self-healing extension (Appendix~\ref{sec:selfheal}) shows how the
union output can serve as a nonce source to recover post-compromise security
with zero additional communication.

\textbf{IBLT techniques.} Goodrich and Mitzenmacher~\cite{GM11} introduced
IBLTs. Rateless IBLTs~\cite{YGA24} (SIGCOMM'24) eliminate a-priori threshold
selection, avoiding resize overhead. Fleischhacker et al.~\cite{FLS23}
(Eurocrypt'23) used IBLTs for compressing encrypted data. Our contribution is
exploiting IBLT \emph{linearity} for incremental delta maintenance, which is
orthogonal to these directions.

% ============================================================
\section{Discussion and Conclusion}
% ============================================================

We presented F-PSU, the first forward-private updatable PSU protocol. F-PSU
achieves per-epoch communication and computation $O(|\Delta|)$, with
$643\times$ communication reduction at $0.1\%$ update ratio, formal UC
security against semi-honest adaptive adversaries, and a provably optimal
leakage profile. The protocol is built from three clean ideas: delta IBLT
encoding exploiting IBLT linearity, split-path incremental peeling with a
provable cascade bound, and a PRF-based key chain for forward privacy with
negligible overhead.

\smallskip\noindent\textbf{Significance.}
F-PSU opens the study of updatable PSU. Prior to this work, every PSU protocol
required re-encoding the full set on each computation. The delta IBLT mechanism
shows that IBLT linearity can be leveraged not just for static union
peelability~\cite{PT26}, but for incremental state maintenance across epochs.
This is a conceptual shift: IBLTs become persistent data structures rather than
ephemeral encoding artifacts.

The forward privacy guarantee addresses a practical concern absent from static
protocols. In any deployment where keys protect sensitive historical data, key
compromise retroactively exposes the entire history unless forward-secure key
management is in place. F-PSU's PRF chain provides this guarantee with zero
additional communication and $<1$~ms overhead, making it practical for
real-world deployments subject to data retention regulations.

\smallskip\noindent\textbf{On the leakage lower bound.}
Theorem~\ref{thm:lower} establishes that any $1$-updatable PSU protocol must
leak individual update counts. This is not a protocol flaw but an inherent
efficiency-leakage trade-off: sublinear communication forces message sizes to
encode information about update magnitude. F-PSU accepts this inevitable
leakage and ensures nothing else is revealed. The bound generalizes beyond
PSU to any two-party stateful computation with incremental updates and
sublinear communication, suggesting a broader theoretical framework for
updatable secure computation.

\smallskip\noindent\textbf{Limitations and future work.}
\emph{Semi-honest security}: the current protocol assumes honest-but-curious
behavior. A malicious $P_0$ could, for instance, apply inconsistent delta IBLT
updates to probe $P_1$'s set. Upgrading to malicious security requires
verifiable delta IBLT integrity---e.g., via commitments to the pre-update IBLT
state combined with zero-knowledge proofs of correct delta application.
Pu et al.~\cite{PuGT26} provide a maliciously secure static IBLT-based PSU,
offering a starting point for the per-epoch component.

\emph{Two-party only}: extending to $m > 2$ parties is natural given IBLT
additivity---$m$ parties can maintain IBLTs whose virtual sum encodes the
multi-party union. The challenge is efficient multi-party UnionPeel, for which
Dong+25~\cite{Dong+25}, GNT24~\cite{GNT24}, and circuit-based PSU~\cite{CGS24} provide building blocks for this extension.

\emph{Post-compromise security}: the self-healing extension
(Appendix~\ref{sec:selfheal}) leverages the union output as an entropy source,
healing in one epoch after a non-trivial honest update. This is a novel
cryptographic use of PSU's inherent output entropy---the union itself becomes
a nonce---and warrants deeper analysis in the multi-party and malicious
settings.

\emph{Post-quantum security}: Gold OPRF~\cite{YBH+25} and LPN-based silent
VOLE~\cite{RRT23} suggest that the OT/OPRF substrate can be replaced with
post-quantum alternatives without changing F-PSU's architecture.

\emph{Performance engineering}: multi-threaded peeling (parallel per-cell
OT/OPRF), hardware-accelerated ChaCha20, and batched OT amortization are
natural optimizations. Rateless IBLTs~\cite{YGA24} would eliminate a-priori
capacity tuning and the occasional full re-encode triggered by set size
growth.

\smallskip\noindent\textbf{Outlook.}
F-PSU establishes that IBLT-based PSU protocols can be made updatable with
provable forward privacy and optimal leakage, while retaining the practical
efficiency that made static IBLT-based PSU competitive with PSI
in the first place~\cite{PT26}. The three core ideas---delta IBLT encoding,
split-path incremental peeling, and PRF-based key evolution---are
conceptually simple and compose cleanly with existing and future PSU
optimizations. We believe updatable PSU will become a standard primitive in
the secure computation toolbox, much as updatable PSI has over the past
five years, and that F-PSU provides a solid foundation for this line of
work. All code is open source~\cite{code}.

% ============================================================
\section{Ethical Considerations}
% ============================================================

\textbf{Stakeholder analysis.} This work develops a privacy-enhancing
cryptographic protocol. Primary stakeholders: end users whose private data is
protected; organizations deploying PSU for collaborative computation; the
research community. The leakage lower bound (Theorem~\ref{thm:lower})
delineates inherent limits of updatable PSU, preventing overclaiming of
privacy properties.

\textbf{Data and human subjects.} All experiments use synthetic data. No human
subjects data, real user information, or production telemetry was used.

\textbf{Vulnerability disclosure.} No vulnerabilities in existing systems are
reported. This work does not break or weaken any deployed cryptographic system.

\smallskip
\noindent\textit{Per USENIX Security policy, the authors attest that:}
\begin{enumerate}
\item This paper presents a novel cryptographic protocol for privacy
enhancement; to the best of our knowledge, it raises no ethical issues beyond
those common to cryptographic research.
\item The authors have complied with the USENIX Security ethical guidelines.
\item No external ethical review was sought, as the work involves only
synthetic benchmark data and no human subjects.
\end{enumerate}

% ============================================================
\section{Open Science}
% ============================================================

\textbf{Artifact availability.} The complete C++ implementation of F-PSU is
available at~\cite{code}. The repository includes build instructions and
scripts to reproduce all experiments in Section~\ref{sec:eval}.

\textbf{Reproducibility.} All benchmarks on Apple M5 (macOS, single-threaded)
using libOTe v2.1. Network conditions are simulated in-app via configurable
RTT and bandwidth caps, reproducible on any modern macOS system. IBLT
parameters ($k = 4$, $e = 3.0$, $\sigma = 64$) and benchmark parameters are
fully documented.

\textbf{Pre-registration.} This work was not pre-registered.

\smallskip
\noindent\textit{Per USENIX Security policy, the authors attest that:}
\begin{enumerate}
\item The open-source code is available at the repository in~\cite{code}.
\item All experimental data can be reproduced using the provided code and
documented parameters.
\item The digital object identifier for the artifact is to be assigned.
\end{enumerate}

% ============================================================
\begin{thebibliography}{99}

\bibitem{KRT19} V.~Kolesnikov, M.~Rosulek, N.~Trieu, X.~Wang.
\newblock Scalable Private Set Union from Symmetric-Key Techniques. \emph{Asiacrypt}, 2019.

\bibitem{JLZ20} Y.~Jia, S.-F.~Sun, H.-S.~Zhou, J.~Du, D.~Gu.
\newblock Shuffle-based Private Set Union: Faster and More Secure. \emph{USENIX Security}, 2022.

\bibitem{ZCL+23} C.~Zhang, Y.~Chen, W.~Liu, M.~Zhang, D.~Lin.
\newblock Linear Private Set Union from Multi-Query Reverse Private Membership Test. \emph{USENIX Security}, 2023.

\bibitem{Tu+25} B.~Tu, Y.~Bai, C.~Zhang, Y.~Cao, Y.~Chen.
\newblock Fast Enhanced Private Set Union in the Balanced and Unbalanced Scenarios. \emph{USENIX Security}, 2025.

\bibitem{DCW13} C.~Dong, L.~Chen, Z.~Wen.
\newblock When Private Set Intersection Meets Big Data: An Efficient and Scalable Protocol. \emph{CCS}, 2013.

\bibitem{IKN+20} M.~Ion, B.~Kreuter, E.~Nergiz, S.~Patel, S.~Saxena, K.~Seth,
M.~Raykova, D.~Shanahan, M.~Yung.
\newblock On Deploying Secure Computing: Private Intersection-Sum-with-Cardinality. \emph{IEEE EuroS\&P}, 2020.

\bibitem{KRS+19} D.~Kales, C.~Rechberger, T.~Schneider, M.~Senker, C.~Weinert.
\newblock Mobile Private Contact Discovery at Scale. \emph{USENIX Security}, 2019.

\bibitem{PT26} L.~Piske, N.~Trieu.
\newblock Is PSI Really Faster Than PSU? Achieving Efficient PSU with Invertible Bloom Filters. \emph{Eurocrypt}, 2026.

\bibitem{Signal} M.~Marlinspike, T.~Perrin.
\newblock The Double Ratchet Algorithm. \emph{Signal Technical Documentation}, 2016.

\bibitem{BMX22} S.~Badrinarayanan, P.~Miao, T.~Xie.
\newblock Updatable Private Set Intersection. \emph{PoPETS}, 2022.

\bibitem{BMS+24} S.~Badrinarayanan, P.~Miao, X.~Shi, M.~Tromanhauser, R.~Zeng.
\newblock Updatable Private Set Intersection Revisited. \emph{Asiacrypt}, 2024.

\bibitem{LMR+26} J.~Liu, P.~Miao, M.~Rosulek, X.~Shi, J.~Wang.
\newblock Updatable PSI from Symmetric-Key Techniques. \emph{Eurocrypt}, 2026.

\bibitem{ABF+25} F.~Alborch, T.~Chauvier, A.~Faonio, A.~Fontaine, F.~Karako\c{c},
A.~K\"{u}p\c{c}\"{u}, C.~Malek, M.~\"{O}nen.
\newblock Updatable Private Set Intersection and Beyond. \emph{ePrint} 2025/2147.

\bibitem{LTS+26} G.~Ling, P.~Tang, S.-F.~Sun, W.~Qiu.
\newblock Efficient Updatable PSI From Asymmetric PSI and PSU. \emph{IEEE TIFS}, 2026.

\bibitem{WZC+24} R.~Wang, J.~Zhou, Z.~Cao, X.~Dong, K.-K.~R.~Choo.
\newblock Updatable PSI with Forward Privacy. \emph{IEEE TIFS}, 2024.

\bibitem{Can01} R.~Canetti.
\newblock Universally Composable Security. \emph{FOCS}, 2001.

\bibitem{libOTe} P.~Rindal et al.
\newblock libOTe: Fast, Portable Oblivious Transfer Library. \url{https://github.com/osu-crypto/libOTe}.

\bibitem{GM11} M.~Goodrich, M.~Mitzenmacher.
\newblock Invertible Bloom Lookup Tables. \emph{Allerton Conference}, 2011.

\bibitem{Ber08} D.~J.~Bernstein.
\newblock ChaCha, a variant of Salsa20. \emph{Workshop Record of SASC}, 2008.

\bibitem{BCK96} M.~Bellare, R.~Canetti, H.~Krawczyk.
\newblock Keying Hash Functions for Message Authentication. \emph{CRYPTO}, 1996.

\bibitem{RS21} P.~Rindal, P.~Schoppmann.
\newblock VOLE-PSI: Fast OPRF and Circuit-PSI from Vector-OLE. \emph{Eurocrypt}, 2021.

\bibitem{RR22} S.~Raghuraman, P.~Rindal.
\newblock Blazing Fast PSI from Improved OKVS and Subfield VOLE. \emph{CCS}, 2022.

\bibitem{RRT23} S.~Raghuraman, P.~Rindal, T.~Tanguy.
\newblock Expand-Convolute Codes for Pseudorandom Correlation Generators from LPN. \emph{CRYPTO}, 2023.

\bibitem{KKRT16} V.~Kolesnikov, R.~Kumaresan, M.~Rosulek, N.~Trieu.
\newblock Efficient Batched Oblivious PRF with Applications to Private Set Intersection. \emph{CCS}, 2016.

\bibitem{PuGT26} S.~Pu, J.~Gao, N.~Trieu.
\newblock Malicious Private Set Union with Two-Sided Output. \emph{Eurocrypt}, 2026.

\bibitem{Bos16} R.~Bost.
\newblock $\Sigma o\varphi o\varsigma$: Forward Secure Searchable Encryption. \emph{CCS}, 2016.

\bibitem{Mat24} J.~P.~Mattsson.
\newblock Security of Symmetric Ratchets and Key Chains -- Implications for Protocols like TLS 1.3, Signal, and PQ3. \emph{ePrint} 2024/220.

\bibitem{HKL+24} K.~Han, S.~Kim, B.~Lee, Y.~Son.
\newblock Revisiting OKVS-based OPRF and PSI. \emph{Asiacrypt}, 2024.

\bibitem{code} Anonymous.
\newblock F-PSU Prototype Implementation. \url{https://github.com/} (released upon publication).

\bibitem{Jia+24} Y.~Jia, S.-F.~Sun, H.-S.~Zhou, D.~Gu.
\newblock Scalable Private Set Union, with Stronger Security. \emph{USENIX Security}, 2024.

\bibitem{Dong+25} M.~Dong, C.~Zhang, Y.~Bai, Y.~Chen.
\newblock Efficient Multi-Party Private Set Union Without Non-Collusion Assumptions.
\emph{USENIX Security}, 2025.

\bibitem{LBL25} Q.~Liu, J.~Bae, J.~Lee.
\newblock Leakage-Free Enhanced Private Set Union. \emph{ePrint} 2025/2087.

\bibitem{YGA24} L.~Yang, Y.~Gilad, M.~Alizadeh.
\newblock Practical Rateless Set Reconciliation. \emph{SIGCOMM}, 2024.

\bibitem{FLS23} N.~Fleischhacker, K.~G.~Larsen, M.~Simkin.
\newblock How to Compress Encrypted Data. \emph{Eurocrypt}, 2023.

\bibitem{GNT24} J.~Gao, S.~Nguyen, N.~Trieu.
\newblock Toward a Practical Multi-Party Private Set Union. \emph{PoPETs}, 2024.

\bibitem{CGS24} G.~R.~Chandran, T.~Schneider, M.~Stillger, C.~Weinert.
\newblock Concretely Efficient PSU via Circuit PSI. \emph{ASIACCS}, 2025.

\bibitem{YBH+25} Y.~Yang, F.~Benhamouda, S.~Halevi, H.~Krawczyk, T.~Rabin.
\newblock Gold OPRF: Post-Quantum Oblivious Power-Residue PRF. \emph{IEEE S\&P}, 2025.

\end{thebibliography}

\newpage
\appendix

% ============================================================
\section{Leakage Lower Bound: Full Proof}
\label{sec:lower-proof}

\begin{proof}[Full proof of Theorem~\ref{thm:lower}]
Fix $P_1$'s set $Y$ and deltas $\Delta^+ Y = \Delta^- Y = \Delta^- X =
\emptyset$. Let $P_0$ add a set $S$ of size $k \in [0, N]$. For each $k$,
let $C(k)$ be the total communication when $|S| = k$.

\textbf{Claim 1.} If $\exists k \neq k'$ with $C(k) \neq C(k')$, then
$|\Delta^+ X|$ is leaked. \emph{Proof:} The semi-honest adversary computes
$|\mathsf{trans}| = C(|S|)$. Since $C$ is non-constant, the adversary learns
$|S|$.

\textbf{Claim 2.} If $C(k) = C_{\max}$ for all $k$, then $C_{\max} \geq
N(\sigma - \log N)$. \emph{Proof:} $S$ a random $N$-subset of $U$. After the
protocol, $P_1$ outputs $S$ (since $Y = \emptyset$). The transcript must
encode $S$, which has entropy $\log\binom{2^\sigma}{N} \geq N(\sigma - \log
N)$ bits. By Shannon's source coding theorem, $C_{\max} \geq H(S)$.

\textbf{Claim 3.} If $\Pi$ is $1$-updatable, $C(k)$ cannot be constant.
\emph{Proof:} If $C(k) = C_{\max}$ for all $k$, then $C(1) = \Omega(N\sigma)$,
contradicting $1$-updatability ($C(k) = O(k \cdot \poly(\sigma, \kappa))$).

The claims imply the theorem for $|\Delta^+ X|$. Symmetric arguments cover
$|\Delta^+ Y|$, $|\Delta^- X|$, $|\Delta^- Y|$.
\end{proof}

% ============================================================
\section{Extended Protocol Analysis and Implementation}
\label{sec:ext-protocol-main}
% ============================================================

\subsection{Design Decisions and Alternatives}

We document design alternatives considered during the development of F-PSU
and the rationale for the final design choices.

\smallskip\noindent\textbf{Why deferred addition merge?}
A natural alternative is to apply both deletions and additions to the main
IBLT in Round~1, then run a single UnionPeel on the modified cells. This
fails because addition cells and deletion cells interact: deleting element~$a$
and adding element~$b$ that share hash cell~$(i,j)$ creates a non-singleton
whose sum is $b - a$, which is neither $a$ nor $b$. The peel thus fails to
recover either element. Separating additions into their own compact IBLT
avoids this cross-contamination; the addition IBLT is fresh (no prior peeling
state), so its peel is clean and deterministic.

\smallskip\noindent\textbf{Why not a single UnionPeel on the delta IBLT?}
Encoding the entire delta $\Delta$ into one compact IBLT and running
UnionPeel on it avoids the split-path complexity. However, the delta IBLT
encodes both additions \emph{and} the \emph{negation} of deletions
($\Encode(\Delta^+) \ominus \Encode(\Delta^-)$), producing a mixed-sign
IBLT. While IBLT subtraction is well-defined, the union-peel analysis of
Piske and Trieu~\cite{PT26} assumes non-negative counts; mixed-sign IBLTs
can create phantom singletons (cells where count~$=1$ but the sum does not
correspond to any element in $X \cup Y$) and false negatives (missed elements
due to cancellation). The split-path design avoids these pathologies by
keeping the addition and deletion paths algebraically clean.

\smallskip\noindent\textbf{Why PRF chain over DH ratchet?}
Signal's double ratchet~\cite{Signal} achieves both forward secrecy and
post-compromise security via ephemeral Diffie-Hellman key exchange. We choose
the PRF chain for F-PSU's main construction because: (i)~it adds zero
communication (DH requires an extra round per epoch); (ii)~it adds negligible
computation (one HMAC-SHA256 vs.\ one X25519 scalar multiplication);
(iii)~forward privacy alone is sufficient for many deployment scenarios where
post-compromise recovery is handled out-of-band. The self-healing extension
(Appendix~\ref{sec:selfheal}) provides post-compromise security when needed,
using only the union output entropy.

\smallskip\noindent\textbf{Why not encrypt the IBLT with a fresh random key?}
Generating a fresh random key per epoch and discarding the old one also
provides forward privacy without a PRF chain. However, this requires the key
to be stored alongside the ciphertext (or derived from external entropy),
increasing storage by $T \cdot \kappa$ bits over $T$ epochs. The PRF chain
requires storing only the current key (256 bits), with all past keys
recoverable if-and-only-if the current key is known---which is exactly the
forward privacy guarantee: past keys are unrecoverable after erasure.

\subsection{Detailed Complexity Analysis}

We provide a detailed per-round breakdown of the communication and
computation costs for epoch $t$.

\smallskip\noindent\textbf{Round~1 (Local IBLT Update).}
\begin{itemize}
\item \emph{Decryption:} $O(m)$ ChaCha20 blocks, where $m = k\ell$ is the
IBLT cell count. Hardware-accelerated on modern processors (AES-NI/ARM
Crypto Extensions), throughput exceeds 10~GB/s, making this negligible
($<1$~ms for $n = 2^{20}$ at $e = 3.0$).
\item \emph{Deletion application:} $k \cdot |\Delta^-|$ cell-wise subtraction
operations (modular arithmetic in $\Z_M$ plus integer decrement). Zero
communication.
\item \emph{Addition encoding:} $k \cdot |\Delta^+|$ cell-wise addition
operations to build $\mathsf{IBLT}_{\Delta}$. Zero communication.
\item \emph{Total computation:} $O(m + k|\Delta|)$ local operations. The
$O(m)$ decrypt/re-encrypt is the dominant term for $|\Delta| \ll m$.
\end{itemize}

\smallskip\noindent\textbf{Round~2a (Mini UnionPeel).}
\begin{itemize}
\item \emph{Communication:} $O(c_k \cdot k \cdot |\Delta^+| \cdot (\sigma +
\kappa))$ bits of OT extension and OPRF traffic. Each of the
$m_{\Delta} = c_k \cdot k \cdot |\Delta^+|$ cells participates in one OT
(for sum retrieval) and one OPRF evaluation (for tag comparison). The total
is deterministic (no cascade) because the addition IBLT has no prior peeling
state.
\item \emph{Computation:} $O(m_{\Delta})$ hash evaluations (IBLT hash
functions), $O(m_{\Delta})$ OT extension operations, and $O(m_{\Delta})$
OPRF evaluations.
\end{itemize}

\smallskip\noindent\textbf{Round~2b (Main UnionPeel).}
\begin{itemize}
\item \emph{Communication:} $O(k \cdot |\Delta^-| \cdot (1 + C_0) \cdot
(\sigma + \kappa))$ bits, where $C_0 = O(k^2|\Delta^-|/m) \ll 1$ bounds
the cascade amplification. By Lemma~\ref{lem:cascade}, total cells examined
$\leq 2k|\Delta^-|$ with overwhelming probability.
\item \emph{Computation:} $O(k|\Delta^-| \cdot (1 + C_0))$ OT and OPRF
operations, plus $O(k|\Delta^-|)$ local deletions.
\end{itemize}

\smallskip\noindent\textbf{Round~3 (Key Evolution).}
\begin{itemize}
\item \emph{Communication:} Zero.
\item \emph{Computation:} One PRF evaluation (HMAC-SHA256, $<1\mu$s) plus
one ChaCha20 encryption of the IBLT ($O(m)$ stream cipher, hardware
accelerated). Secure erasure of $K_{t-1}$ is a constant-time memory
overwrite.
\end{itemize}

\smallskip\noindent\textbf{Total per-epoch cost.}
Communication: $O(|\Delta| \cdot (\sigma + \kappa) + |X_{\cup}| \cdot
\sigma)$ bits. The $O(|\Delta|)$ term dominates the secure computation
phase; the $O(|X_{\cup}|)$ term is the inherent PSU output (transmitting the
peeled union elements from $P_1$ to $P_0$). Computation: $O(m + k|\Delta|)$
local operations plus $O(k|\Delta|)$ OT/OPRF operations.

\subsection{IBLT Parameter Selection}

The IBLT expansion factor $e = m/(k \cdot \max(|X|, |Y|))$ and hash count
$k$ are chosen to satisfy Theorem~\ref{thm:iblt}. We use $k = 4$ and
$e = 3.0$ throughout, giving $m = 4 \cdot 3.0 \cdot \max(|X|, |Y|)$ cells.
With these parameters, the $2$-core emptiness probability exceeds
$1 - 2^{-40}$ for $|X| + |Y| \leq 2^{20}$, sufficient for all our
experiments.

The expansion $e = 3.0$ is conservative: Theorem~\ref{thm:iblt} only requires
$e > c_4 \approx 1.295$. The extra margin accommodates set size growth between
full re-encodes (see Section~\ref{sec:eval}, Incremental Peeling and Parameter
Growth) and provides headroom for the cascade bound.

\smallskip\noindent\textbf{Resize policy.}
When $|X_t| + |Y_t|$ exceeds the IBLT design threshold, a full re-encode is
triggered: both parties encode their current sets into a larger IBLT and run
a fresh UnionPeel. The resize cost is $O(n)$ but amortized over
$\Omega(n/\text{growth\_rate})$ epochs. With exponential growth (doubling
capacity each resize) and linear set growth ($n_t = n_0 + t \cdot g$), the
amortized per-epoch resize cost is $O(g)$, matching the delta cost.

% ============================================================
\section{Extended Experimental Results}
\label{sec:ext-eval}
% ============================================================

We provide additional experimental data complementing
Section~\ref{sec:eval}.

\subsection{Parameter Sensitivity}

We evaluate epoch time sensitivity to hash count $k$ and expansion factor
$e$ at $n = 2^{16}$, $|\Delta| = 2^{10}$. Increasing $k$ improves peeling success probability but linearly increases
encode and peel costs. The expansion $e = 3.0$ with $k = 4$ provides a
balanced trade-off: strong probabilistic guarantees ($\Pr[\text{failure}]
< 2^{-40}$) with moderate cell count overhead.

\subsection{Memory Footprint}

Each IBLT cell stores a 64-bit sum ($\Z_{2^{64}}$) and a 16-bit count,
for 10 bytes per cell. With $k = 4$, $e = 3.0$, $n = 2^{20}$: $m =
4 \cdot 3.0 \cdot 2^{20} \approx 12.6$~million cells, occupying
$\approx 126$~MB. The compact addition IBLT at $0.1\%$ update ratio uses
$m_{\Delta} \approx 12{,}600$ cells ($\approx 126$~KB), a $1000\times$
reduction. The ciphertext size equals the plaintext size (ChaCha20 is a
stream cipher).



\subsection{Microbenchmarks: Extended}

We benchmark IBLT serialization (for TCP transmission) and various OT
extension configurations. At $n = 2^{16}$ with $e = 3.0$, IBLT
serialization completes in under $2$~ms and OT extension setup adds
under $5$~ms, negligible relative to network latency in all profiles.



\subsection{Delta Ratio Sweep}

\begin{figure}[t]
\centering
\includegraphics[width=0.8\columnwidth]{fig_delta_sweep.pdf}
\caption{Epoch time vs.\ delta ratio ($n = 2^{16}$, LAN). Break-even occurs at
$|\Delta|/n \approx 50\%$; below this threshold, F-PSU provides savings.}
\label{fig:delta}
\end{figure}

Fig.~\ref{fig:delta} shows epoch time as a function of the delta ratio
$|\Delta|/n$, at fixed $n = 2^{16}$ under LAN. The break-even point where
delta IBLT outperforms full re-encoding occurs at $|\Delta|/n \approx 50\%$;
below this threshold, F-PSU provides communication savings.



% ============================================================
\section{Full UC Security Proofs}
\label{sec:full-proofs}

\subsection{Full UC Security Proof}
\label{sec:full-proof}
% ============================================================

This appendix provides the complete UC security proof referenced in
Section~\ref{sec:security}. We give the full simulator construction, prove
all hybrid transitions, and establish the composition theorem.

\subsection{Complete Simulator Construction}

\begin{framed}
\noindent\textbf{Simulator $\cS$ for $\Pi_{\mathsf{fpsu}}$}

\smallskip\noindent\textbf{Input:} Access to $\cF_{\mathsf{fpsu}}$ and
$\cA$ (corrupting $P_0$ at epoch $\tau$).

\smallskip\noindent\textbf{Internal State:}
\begin{itemize}
\item $\tilde{Y}_t$: simulated $P_1$ set for each epoch $t$
\item $\tilde{K}^S_t \in \{0,1\}^\kappa$: simulated encryption key
\item $\widetilde{\mathsf{CT}}^{(t)}_S$: simulated ciphertext
\item $\tilde{\sigma}_{\mathsf{VOLE}}$, $\tilde{\sigma}_{\mathsf{OKVS}}$:
simulated VOLE and OKVS states
\end{itemize}

\smallskip\noindent\textbf{Procedure:}

\smallskip\noindent\textit{Initialization} ($t = 0$):
\begin{enumerate}
\item Sample $\tilde{K}^S_0, \tilde{K}^C_0 \sample \{0,1\}^\kappa$.
\item Initialize empty IBLTs for both parties.
\item Simulate VOLE setup: $\tilde{\sigma}_{\mathsf{VOLE}} \leftarrow
\textsf{SimVOLE}(1^\kappa)$. In the $\cF_{\mathsf{VOLE}}$-hybrid model,
$\cS$ simply records the ideal VOLE handles.
\end{enumerate}

\smallskip\noindent\textit{Per-Epoch Simulation} (epoch $t < \tau$):
\begin{enumerate}
\item Receive $X_{\cup,t}$ from $\cF_{\mathsf{fpsu}}$; receive $X_t$ (as
$P_0$'s input) from $\cZ$ through $\cA$.
\item Compute $\tilde{Y}_t \leftarrow X_{\cup,t} \setminus X_t$. By
construction, $X_t \cup \tilde{Y}_t = X_{\cup,t}$.
\item \textbf{Simulate UnionPeel.} Compute the virtual union IBLT:
$\mathsf{IBLT}_\cup \leftarrow \Encode(X_{\cup,t})$. Run the full
$\List(\mathsf{IBLT}_\cup)$ algorithm to obtain the peeling trace---the
ordered sequence of peeled elements and their originating cells
$(v_1, (i_1, j_1)), (v_2, (i_2, j_2)), \dots$

\item \textbf{Program ideal OT.} For each cell $(i,j)$ in the query set
$Q^{(r)}$ at iteration $r$, the ideal $\cF_{\mathsf{OT}}$ receives $P_0$'s
choice bit $c_{i,j}$ and $P_1$'s messages $(m_0, m_1)$. $\cS$ programs the
output to match the peeling trace:
\begin{itemize}
\item If the peeled element $v_{i,j}$ is unique to $P_0$
($\mathsf{cnt}_C = 1, \mathsf{cnt}_S = 0$): the 1-out-of-2 OT returns
$f_{i,j} \leftarrow m_{c_{i,j}}$ such that the subsequent computation
yields $v_{i,j} = \mathsf{sum}_C[i][j]$.
\item If $v_{i,j}$ is unique to $P_1$
($\mathsf{cnt}_C = 0, \mathsf{cnt}_S = 1$): the OT returns
$f_{i,j} \leftarrow \mathsf{sum}_S[i][j]$ from $P_1$'s message.
\item If $v_{i,j}$ is shared ($\mathsf{cnt}_C = \mathsf{cnt}_S = 1,
\mathsf{sum}_C = \mathsf{sum}_S$): the 1-out-of-3 OT result selection
returns $\mathsf{sum}_C[i][j]$.
\end{itemize}
$\cS$ records the programmed outputs in its local OT transcript store.

\item \textbf{Program ideal OPRF.} For each cell $(i,j)$, $\cS$ receives
$P_1$'s OPRF query $z_{i,j}$ and must produce tag $y_{i,j}$. If
$\mathsf{cnt}_C = 1$ and $\mathsf{sum}_C = z_{i,j}$ (equality case), $\cS$
programs $y_{i,j} \leftarrow F_k(\mathsf{sum}_C[i][j])$ using the known
OPRF key. Otherwise, $\cS$ outputs a uniformly random tag. This
programming is always possible because $\cS$ knows the full
$\mathsf{IBLT}_\cup$ and can precompute the equality conditions.

\item \textbf{Simulate ciphertext.} Sample $\tilde{K}^S_t \sample
\{0,1\}^\kappa$ (independent random). Encrypt a dummy IBLT of the
appropriate size: $\widetilde{\mathsf{CT}}^{(t)}_S \leftarrow
\ChaCha.\Enc(\tilde{K}^S_t, N_S, 0^{|\mathsf{IBLT}_S|})$.

\item Simulate erasure: overwrite $\tilde{K}^S_{t-1}$ with zeros.
\end{enumerate}

\smallskip\noindent\textit{Corruption Handling} (epoch $t = \tau$):
\begin{enumerate}
\item $\cF_{\mathsf{fpsu}}$ sends $(X_\tau, K^C_\tau)$ to $\cS$.
\item $\cS$ delivers $(X_\tau, K^C_\tau)$ to $\cA$ as $P_0$'s internal
state.
\item $\cS$ continues simulating $P_1$ for epochs $t \geq \tau$ using the
same per-epoch procedure as above.
\end{enumerate}
\end{framed}

\subsection{Hybrid Argument: Full Proofs}

We prove Theorem~\ref{thm:main} via three hybrid simulators.

\begin{lemma}[Random Keys]\label{lem:hyb1-full}
$|\Pr[\cZ(\cS_1)=1] - \Pr[\cZ(\cS_0)=1]| \leq T \cdot
\mathsf{Adv}^{\mathsf{PRF}}_{F}(\kappa)$.
\end{lemma}

\begin{proof}
We construct a reduction $\cB$ to the PRF security of
$F = \mathsf{HMAC\text{-}SHA256}$. $\cB$ receives oracle access to either
$F(K, \cdot)$ (real PRF) or a truly random function $R(\cdot)$.

$\cB$ simulates the protocol execution for $\cZ$ using the hybrid
construction $\cS_1$, with the following modification for epoch $t$:
instead of sampling $\tilde{K}^S_t \sample \{0,1\}^\kappa$, it queries the
oracle on the fixed domain separator $c$: $\tilde{K}^S_t \leftarrow
\mathcal{O}(c)$.

If $\mathcal{O} = F(K, \cdot)$, then $\tilde{K}^S_t = F(K, c)$.
Combined with the real key chain starting from $K^S_0 = K$, this produces
exactly the key distribution of $\cS_0$ (real protocol).

If $\mathcal{O} = R(\cdot)$, then $\tilde{K}^S_t = R(c)$ is independent
of $K$ and of all other keys, producing exactly the distribution of
$\cS_1$ (random keys).

$\cB$ outputs whatever $\cZ$ outputs. The distinguishing advantage of
$\cB$ is exactly $|\Pr[\cZ(\cS_1)=1] - \Pr[\cZ(\cS_0)=1]|$, which is
bounded by $\mathsf{Adv}^{\mathsf{PRF}}_{F}(\kappa)$ for a single epoch.
Summing over $T$ epochs via a standard hybrid argument yields the $T
\cdot \mathsf{Adv}^{\mathsf{PRF}}_{F}(\kappa)$ bound.
\end{proof}

\begin{lemma}[Dummy Ciphertexts]\label{lem:hyb2-full}
$|\Pr[\cZ(\cS_{1.5})=1] - \Pr[\cZ(\cS_1)=1]| \leq T \cdot
\mathsf{Adv}^{\mathsf{PRG}}_{\mathsf{ChaCha20}}(\kappa)$.
\end{lemma}

\begin{proof}
Analogous to Lemma~\ref{lem:hyb1-full}. For each ciphertext, replace the
ChaCha20 keystream $\ChaCha(\tilde{K}^S_t, N_S)$ with a truly random
string of the same length. Under the PRG security of ChaCha20, the
distributions are indistinguishable with advantage $\leq
\mathsf{Adv}^{\mathsf{PRG}}(\kappa)$ per ciphertext, and $\leq T \cdot
\mathsf{Adv}^{\mathsf{PRG}}(\kappa)$ over $T$ epochs. Note that
$\tilde{K}^S_t$ is already independent random (from Lemma~\ref{lem:hyb1-full}),
so the PRG reduction is straightforward.
\end{proof}

\begin{lemma}[Simulated Inputs]\label{lem:hyb3-full}
$\Pr[\cZ(\cS_2)=1] = \Pr[\cZ(\cS_{1.5})=1]$.
\end{lemma}

\begin{proof}
In both $\cS_{1.5}$ and $\cS_2$, $P_1$'s keys are independent random and
the IBLT ciphertexts are encryptions of the all-zero string. The only
difference is the content of $P_1$'s IBLT fed into
$\Pi_{\mathsf{UnionPeel}}$.

By Theorem~\ref{thm:upeel} (Union Peel Equivalence),
$\uPeel(\mathsf{IBLT}_C, \mathsf{IBLT}_S, i, j) =
\Peel(\mathsf{IBLT}_\cup, i, j)$ where $\mathsf{IBLT}_\cup =
\mathsf{IBLT}_C \oplus \mathsf{IBLT}_S$. The virtual union depends only
on $X_t \cup Y_t$, which equals $X_t \cup \tilde{Y}_t = X_{\cup,t}$ by
construction of $\tilde{Y}_t$. Therefore the peeling outputs---and hence
the entire protocol transcript in the
$(\cF_{\mathsf{OPRF}}, \cF_{\mathsf{OT}})$-hybrid model---are identical
in distribution for $Y_t$ and $\tilde{Y}_t$.

After corruption at epoch $\tau$, $\cF_{\mathsf{fpsu}}$ reveals
$(X_\tau, K^C_\tau)$. The simulated view for $t < \tau$ used only
$X_{\cup,t}$, which $\cS$ received from $\cF_{\mathsf{fpsu}}$. No
inconsistency arises because $\cS$ never committed to a specific $Y_t$
beyond what $X_{\cup,t}$ implies. The view is perfectly identical.
\end{proof}

Combining Lemmas~\ref{lem:hyb1-full}-\ref{lem:hyb3-full}, the total
distinguishing advantage between the real world ($\cS_0$) and the ideal
world ($\cS_2$) is bounded by $T \cdot (\mathsf{Adv}^{\mathsf{PRF}}(\kappa)
+ \mathsf{Adv}^{\mathsf{PRG}}(\kappa)) \leq \negl(\kappa)$, completing the
proof of Theorem~\ref{thm:main}. \hfill$\square$

\subsection{UC Composition: Formal Statement}

\begin{theorem}[UC Composition]\label{thm:composition}
Let $\Pi_{\mathsf{epoch}}$ be a protocol that UC-realizes $\cF_{\mathsf{psu}}$
in the $(\cF_{\mathsf{OPRF}}, \cF_{\mathsf{OT}})$-hybrid model, and let
$\Pi_{\mathsf{key}}$ be a protocol that UC-realizes $\cF_{\mathsf{fwKey}}$.
Then the parallel composition $\Pi_{\mathsf{epoch}} \parallel
\Pi_{\mathsf{key}}$ UC-realizes $\cF_{\mathsf{fpsu}}$ in the same hybrid
model.
\end{theorem}

\begin{proof}[Proof sketch]
By the universal composition theorem~\cite{Can01}, it suffices to exhibit a
simulator for $\cF_{\mathsf{fpsu}}$ given simulators for $\cF_{\mathsf{psu}}$
and $\cF_{\mathsf{fwKey}}$. The construction is modular: the per-epoch
simulator $\cS_{\mathsf{epoch}}$ (from Piske and Trieu~\cite{PT26}) handles
all OT/OPRF interactions within each epoch, receiving $X_{\cup,t}$ from
$\cF_{\mathsf{fpsu}}$. The key management simulator $\cS_{\mathsf{key}}$
samples independent random keys for $t < \tau$ (indistinguishable from PRF
outputs by PRF security) and delivers the current key upon corruption. The
composition is a standard application of the UC theorem, requiring no new
cryptographic assumptions.
\end{proof}

% ============================================================
\section{Self-Healing Key Chains}
\label{sec:selfheal}
% ============================================================

F-PSU achieves forward security but not post-compromise security. We present
a protocol extension achieving \emph{self-healing} with \textbf{zero additional
communication}, protecting against single-party compromise.

\textbf{Key insight: the union output as a nonce source.}
At each epoch, both parties learn $X_{\cup,t}$. Critically, \emph{neither
party knows which elements were contributed by the other}. From $P_0$'s
perspective, $X_{\cup,t} \setminus X_t$ is unpredictable (it consists of
$P_1$'s exclusive elements). This entropy can serve as a fresh random nonce.

\textbf{Protocol modification.}
Replace the key update with:
\[
K^b_t \leftarrow \PRF(K^b_{t-1},\; c \concat H(X_{\cup,t})),
\]
where $H$ is a collision-resistant hash. Both parties compute $H(X_{\cup,t})$
locally---no additional round, no additional communication.

\textbf{Security argument.}
Consider compromise at epoch $\tau$. The adversary learns $K^b_\tau$ and
$X_{\cup,\tau}$. To compute $K^b_{\tau+1}$, the adversary needs
$H(X_{\cup,\tau+1})$ at the next epoch. But $X_{\cup,\tau+1} \setminus
X_{\cup,\tau}$ contains deltas applied by the \emph{honest party}, which the
adversary cannot predict. Therefore:
\begin{itemize}
\item Before compromise: adversary without $X_{\cup,t}$ (forward privacy
prevents decryption) cannot compute $K^b_{t+1}$.
\item After compromise detected: the honest party's delta introduces fresh
entropy, healing the chain.
\item Healing time: one epoch after non-trivial honest update.
\end{itemize}

\textbf{Security bound.}
For exclusive set size $|X_{\cup,t} \setminus X_b| \geq q$, min-entropy:
$\log\binom{|U|-|X_b|}{q} \geq q(\sigma - \log|U|)$ bits. With $q = 1$ and
$\sigma \geq 128$, this exceeds $\kappa = 128$, giving single-epoch healing.

\textbf{Practical considerations.}
Healing requires the honest party to make a non-trivial update. \emph{Entropy
decay}: if the adversary observed $T$ prior unions, the honest party's
exclusive set may be reduced, increasing expected healing time to
$O(|U|/|\text{exclusive}|)$ epochs. This construction shows that PSU's
inherent output entropy can be leveraged for cryptographic self-healing
without additional communication, a property unique to updatable PSU among
secure computation primitives.

\subsection{Formal Security Definition}

We formalize the self-healing property via a game-based definition.

\begin{definition}[SH-IND Game]
Let $\Pi$ be an updatable PSU protocol with self-healing key chains. The
$\mathsf{SH\text{-}IND}$ game with security parameter $\kappa$ and $T$
epochs proceeds as follows:
\begin{enumerate}
\item The challenger samples $b \sample \{0,1\}$ and initializes both
parties with random keys $K^0_0, K^1_0 \sample \{0,1\}^\kappa$.
\item The protocol runs for $T$ epochs. The adversary $\cA$ observes all
transcripts and may adaptively choose its own deltas for the corrupted
party.
\item At epoch $\tau$, $\cA$ corrupts $P_0$, receiving
$(X_\tau, K^0_\tau)$.
\item The protocol continues for epoch $\tau+1$. $P_1$ (honest) provides
a non-trivial delta ($|\Delta_{P_1}| \geq 1$).
\item If $b = 0$: $\cA$ receives the real key $K^0_{\tau+1} \leftarrow
\PRF(K^0_\tau, c \concat H(X_{\cup,\tau+1}))$. If $b = 1$: $\cA$
receives $R \sample \{0,1\}^\kappa$ independent of the protocol state.
\item $\cA$ outputs $b'$ and wins if $b' = b$.
\end{enumerate}
The protocol satisfies self-healing if for every PPT $\cA$:
$\mathsf{Adv}^{\mathsf{SH\text{-}IND}}_\cA(\kappa) = |\Pr[b'=b] -
\tfrac{1}{2}| \leq \negl(\kappa)$.
\end{definition}

\begin{theorem}[Self-Healing Security]\label{thm:selfheal}
Assume $F$ is a secure PRF and $H$ is a random oracle. Then the
self-healing construction satisfies the SH-IND definition with
$\mathsf{Adv}^{\mathsf{SH\text{-}IND}}(\kappa) \leq
\mathsf{Adv}^{\mathsf{PRF}}(\kappa) + 2^{-q(\sigma - \log|U|)}$, where
$q = |X_{\cup,\tau+1} \setminus X_\tau|$ is the honest party's exclusive
contribution at the healing epoch.
\end{theorem}

\begin{proof}[Proof sketch]
By PRF security, we can replace $\PRF(K^0_\tau, \cdot)$ with a random
function $R(\cdot)$, incurring advantage $\mathsf{Adv}^{\mathsf{PRF}}$.
The key $K^0_{\tau+1} = R(c \concat H(X_{\cup,\tau+1}))$ is then
indistinguishable from random unless the adversary can predict
$H(X_{\cup,\tau+1})$. By the random oracle property, this requires
predicting $X_{\cup,\tau+1} \setminus X_\tau$---the honest party's delta
at epoch $\tau+1$. The adversary knows $X_\tau$ (from corruption) and
$X_{\cup,\tau}$ (from the protocol output), but the honest party's delta
introduces fresh entropy of at least $q(\sigma - \log|U|)$ bits.
When this exceeds $\kappa$, the adversary's guessing advantage is
$2^{-\kappa}$, negligible.
\end{proof}

\subsection{Comparison with DH-Based Ratchets}

Unlike DH-based ratchets~\cite{Signal} which require an extra round per
epoch for ephemeral key exchange, the self-healing construction uses only
the union output entropy (zero additional communication). Healing occurs
in a single epoch whenever the honest party makes a non-trivial
update---a condition naturally satisfied in deployments with regular
activity (e.g., daily IoC feeds). The approach is complementary to DH:
both can be combined for defense in depth.

% ============================================================
\section{Additional Material}
\label{sec:additional}
% ============================================================

\subsection{Implementation Details}
\label{sec:impl-details}

The C++ implementation (six modules, $\approx$2,000 lines, excluding
libOTe) is available at~\cite{code}. Build and reproduction instructions
are in the repository README. Experiments use an Apple M5 (macOS),
Clang~16, C++20, \texttt{-O3}; IBLT parameters $k=4$, $e=3.0$,
$\sigma=64$ bits.
% ============================================================

\subsection{Formal Cascade Analysis}
\label{sec:formal-cascade}
% ============================================================

We provide the complete proof of the Cascade Bound Lemma
(Lemma~\ref{lem:cascade}) from Section~\ref{sec:protocol}.

\subsection{Proof of Lemma~\ref{lem:cascade}}

\begin{proof}
Let $\mathsf{IBLT}_\cup^{(t-1)} = \Encode(X_{t-1} \cup Y_{t-1})$ be the
virtual union IBLT after full peeling at epoch $t-1$. Since the peel
completed successfully (by the IBLT parameter choice guaranteeing
$2$-core emptiness with probability $\geq 1 - 2^{-\lambda}$), every
cell in $\mathsf{IBLT}_\cup^{(t-1)}$ has count zero.

At epoch $t$, each party applies its delta to its local IBLT. The
virtual union after delta application is:
\[
\mathsf{IBLT}_\cup^{(t)} = \mathsf{IBLT}_\cup^{(t-1)}
  \oplus \Encode(\Delta^+_X \cup \Delta^+_Y)
  \ominus \Encode(\Delta^-_X \cup \Delta^-_Y)
\]
Since $\mathsf{IBLT}_\cup^{(t-1)} = \mathbf{0}$ (all-zero), the initial
state for peeling is equivalent to $\Encode(\Delta)$ on a fresh IBLT,
where $\Delta = \Delta^+_X \cup \Delta^+_Y \cup \Delta^-_X \cup
\Delta^-_Y$, with deletions encoded as negative contributions.

The peeling process on this fresh IBLT follows the standard $\List$
algorithm. Cells with count~$=1$ are singletons whose sum equals the
unique element; these are peeled and deleted, potentially creating new
singletons (the cascade).

We bound the cascade size. Let $Q_0 = \{(i, h_i(x)) : x \in \Delta^-, i
\in [k]\}$ be the initial query set (size $k|\Delta^-|$). After peeling
all singletons reachable from $Q_0$, the remaining cells form a subgraph
of the random $k$-uniform hypergraph induced by the $|\Delta|$ elements.
The expected number of additional cells examined (the cascade) is bounded
by the size of the mini-$2$-core formed by the delta elements alone.

By standard random hypergraph theory~\cite{GM11}, for $|\Delta| \ll
m/k$, the probability that $|\Delta|$ random edges create a non-empty
$2$-core on $m$ vertices is $O(|\Delta|^{-k+2})$. When the $2$-core is
empty, peeling from any initial set of singletons recovers all elements,
and the cascade size is at most $(1 + C_0) \cdot k|\Delta^-|$ where
\[
C_0 = O\left(\frac{k^3 |\Delta^-|^2}{m}\right)
\]
bounds the expected size of the mini-$2$-core. This follows from the
birthday paradox: the expected number of edge pairs sharing $k$ vertices
is $\binom{k|\Delta^-|}{2} \cdot k/m$.

With the IBLT parameters chosen such that $m > c_k \cdot n$ and
$|\Delta^-| \ll n$, we have $|\Delta^-| \ll m/k$, and $C_0 \ll 1$.
Therefore the total cells examined is $\leq (1 + C_0) \cdot k|\Delta^-|
\leq 2k|\Delta^-|$ with overwhelming probability.

The cascade terminates in $O(\log(k|\Delta^-|))$ iterations because each
iteration removes at least a constant fraction of the remaining
singletons (by the random hypergraph expansion properties below the
$2$-core threshold).
\end{proof}

\subsection{Experimental Validation of the Cascade Bound}

We experimentally validated the cascade bound across the full parameter
range. Table~1 in Section~\ref{sec:eval} shows that the
observed cascade ratio is $\approx 1.95 \times k|\Delta|$ (averaged over
20 trials), well within the $2 \times k|\Delta|$ theoretical bound. The
ratio is stable across all values of $n$ and $|\Delta|$, confirming that
the $2$-core emptiness assumption holds in practice for the chosen IBLT
parameters.

The cascade ratio approaches $1.0 \times k|\Delta|$ as $|\Delta|/n \to
0$, since the probability of hash collisions among delta elements
vanishes with decreasing delta density. At $|\Delta|/n = 0.1\%$, we
observe cascade ratios of $1.01$--$1.05 \times k|\Delta|$, indicating
that the delta elements are essentially isolated in the IBLT hypergraph.

% ============================================================
\subsection{Comparison with Related Updatable Protocols}
\label{sec:comparison-detail}

\textbf{FUPSI~\cite{WZC+24}} (TIFS 2024) achieves updatability and forward
privacy for PSI via keyword PIR ($O(\log n)$ communication) and
SparsePIR-based storage. F-PSU targets PSU with a symmetric-key PRF chain
($O(|\Delta|)$ communication) and IBLT-based storage, exploiting IBLT
linearity for native set operations and deletion.

\textbf{LMR+26~\cite{LMR+26}} (Eurocrypt 2026) introduces adaptive input
selection attacks on UPSI. F-PSU is immune: the union output is a superset
of both inputs, so an adversary cannot distinguish whether $x \in
X_{\cup,t}$ due to addition or prior honest possession. LMR+26 does not
provide forward privacy; F-PSU adds it with zero overhead.

\textbf{ePSU~\cite{Tu+25,LBL25}} (USENIX 2025) eliminates within-execution
leakage, orthogonal to F-PSU's cross-epoch forward privacy. The mechanisms
are complementary: an ePSU made updatable via F-PSU's delta IBLT would
achieve both guarantees; we leave this to future work.

% ============================================================
\subsection{Deployment Considerations}
\label{sec:deployment}

F-PSU is most beneficial when: (i)~sets evolve incrementally with
$|\Delta| \ll n$ (the update ratio is small); (ii)~frequent
recomputation of the union is required (low-latency requirements);
(iii)~sensitive historical data must be protected against key compromise
(regulatory compliance); (iv)~network bandwidth is constrained (WAN or
low-bandwidth links). In these scenarios, F-PSU provides substantial
communication savings ($643\times$ at $0.1\%$ update ratio) and forward
privacy with negligible overhead.

F-PSU provides limited benefit when: (i)~the update ratio is high
($|\Delta|/n > 50\%$, see Fig.~\ref{fig:delta}); (ii)~the
deployment is on a high-bandwidth LAN where IBLT peeling computation
dominates; (iii)~forward privacy is not required and a static protocol
suffices.

\smallskip\noindent\textbf{Scalability.}
For $n > 2^{20}$, rateless IBLTs~\cite{YGA24}, sharding, and
multi-threading compose naturally with F-PSU's linear delta mechanism.

\end{document}
